\section{Results and evidence boundary}

\noindent\textbf{6 September: explicit spectral and umbral maps, and the original family's parameter identifications.}
For the literal common lattice $\mathcal I=\{\sqrt2(a,a,a):a\in D_4\}$,
write $a=a_0+v$, $t^2=|v|^2$, $Q=a_0^2+t^2$ and
$P=a_0(a_0^2-3t^2)$. Its complete cubic image is
\[
 \tau(\mathcal I)=2\sqrt2\bigl(18\mathbb Z\cup(6\mathbb Z+2)
                                      \cup(6\mathbb Z-2)\bigr).
\]
Three explicit integral affine lines give a section of this value map.
The scalar-retaining Albert matrix $X=\mu I+Y(a)$, with cyclic
off-diagonal entries $\sqrt2a$, has invariants and labelled spectrum
\[
\begin{aligned}
 (\operatorname{tr}X,\operatorname{tr}X^2,\det X)
   &=(3\mu,3\mu^2+12Q,\mu^3-6\mu Q+4\sqrt2P),\\
 \operatorname{Spec}Y
   &=(2\sqrt2a_0,-\sqrt2a_0-\sqrt6t,-\sqrt2a_0+\sqrt6t).
\end{aligned}
\]
For $t>0$, the explicit labelled eigenprojector satisfies
$(\Pi_1)_{12}=-1/6-\sqrt3(v/t)/6$; thus the frame and spectrum
recover $a$, not only its norm and cubic. The full discriminant is
$864t^2(3a_0^2-t^2)^2$. The full Cartan family retains the scalar
$\mu$, radius $r$, and angle $\theta$, with $r^2=12Q$ and
$r^3\cos3\theta=24\sqrt3P$ on this subfamily.

For the same marked root system $A_5^4D_4$, the
Cheng--Duncan--Harvey prescription gives
$S^X=(5I_{12}+J)S_6$, where $J_{rs}=\mathbf1_{s=5r\ (12)}$ and
$S_{6,r}=\sum_{k\equiv r\ (12)}kq^{k^2/24}$.
The coefficient map to the retained Fricke function is
\[
 T_6=\frac{\eta(\tau)^5\eta(3\tau)}{\eta(2\tau)\eta(6\tau)^5},
 \qquad f^X=-q\frac{d}{dq}\log T_6,
 \qquad T_6\!\left(-\frac1{6\tau}\right)=\frac{72}{T_6(\tau)}.
\]
All twelve theta indices, the multiplier determinant, the original
glue-group action and its twisted shadows are computed in the proofs.
This is an explicit use and comparison of established umbral machinery,
not a claim to a new moonshine module.

In the original analytic family $\beta_c=\beta^{\rm part}+c$,
$\operatorname{Im}c<-M$, the actual global identifications are exactly
$X_c\simeq X_{c'}\Longleftrightarrow c'-c\in2\mathbb Z$.
The middle lattices agree for every integer shift; the multiplicity-four
filling retains an order-two obstruction to odd shifts. Even shifts have
explicit finite-filling and toric cusp maps and a proved global additive
gluing correction. The resulting parameter quotient is
$w=e^{-\pi ic}$, $0<|w|<e^{-\pi M}$.
This classifies the retained analytic family without using the claimed
identification of its total space with $S^6$.


\paragraph{6 September: cubic arithmetic, integral lattice maps, and the canonical ring.}
For the retained $A_5^4\perp D_4$ Niemeier lattice $N$ and its
original index-six Leech neighbour $\Lambda$, fix the Cayley--Dickson
frame $(1,i,j,k,\ell,i\ell,j\ell,k\ell)$ used in the complete cubic
proof. Put $\tau(q)=\operatorname{Re}((q_1q_2)q_3)$ and
$F=\tau\circ\mathcal L$ on the original Niemeier span,
and $\mathcal V_F(L)=\operatorname{span}_{\mathbb Z}\{F(y):y\in L\}$.
The exact additive groups are
\[
\begin{aligned}
 \mathcal V_F(N)&=\tfrac1{27}\mathbb Z\oplus
       \tfrac{\sqrt2}{54}\mathbb Z\oplus\tfrac{\sqrt3}{18}\mathbb Z,\\
 \mathcal V_F(\Lambda)&=\tfrac1{27}\mathbb Z\oplus
       \tfrac{\sqrt2}{216}\mathbb Z\oplus\tfrac{\sqrt3}{54}\mathbb Z,
 \qquad \mathcal V_F(\Lambda)/\mathcal V_F(N)\cong\mathbb Z/12.
\end{aligned}
\]
The Albert determinant is $\lambda^3-\lambda(y,y)+2F(y)$.
The quotient concerns the displayed value-generated groups in this fixed
frame, not an inclusion $N\subset\Lambda$. The actual residue
$r=\rho(361)$ cycles three Leech neighbours $\Lambda_j=r^j\Lambda$.
Their pairwise intersection is one integral lattice $J$, with
$\det J=81$ and $[\Lambda_j:J]=9$. The fixed fourth isotropic plane
in $J^\vee/J$ gives $M$ with exactly $72$ roots of type $A_2^{12}$,
$M\cap\Lambda_j=J$, and $[M:J]=9$.
\[
 \mathcal V_F(J)=\mathcal V_F(M)
 =\tfrac1{27}\mathbb Z\oplus\tfrac{\sqrt2}{216}\mathbb Z
                         \oplus\tfrac{\sqrt3}{18}\mathbb Z.
\]
The successive value-group inclusions have matrices
$\operatorname{diag}(1,4,1)$ and $\operatorname{diag}(1,1,3)$,
giving quotients $\mathbb Z/4$ and $\mathbb Z/3$.
Both containments and their reverse generator certificates are calculated
from the full cubic polarization, not inferred from lattice indices.

Wilson's construction \cite{esn:Wilson2009}, transported by the explicit
isometry $\iota(x)=x/\sqrt2$, gives $\Lambda_{\rm cyc}$ in the
original triality metric. Its cyclic fixed Gram matrix is $3G_{E_8}$;
its orthogonal complement is explicitly isometric to $A_2\otimes E_8$,
and their sum has index $3^8$, with quotient $\mathbb F_3^8$.
The directions remain $c(q_1,q_2,q_3)=(q_3,q_1,q_2)$ and
$c_{\rm W}=c^{-1}$. Its literal intersections are
\[
 \Lambda_{\rm cyc}\cap\mathcal L(N)
 =\Lambda_{\rm cyc}\cap\mathcal L(\Lambda)
 =\Delta:=\{\sqrt2(a,a,a):a\in D_{\mathbb H}\},
\]
where $D_{\mathbb H}$ is the even-coordinate-sum integer lattice in
the first four stated Cayley coordinates. Thus $\Delta$ has rank $4$,
Gram matrix $6G_{D_4}$, determinant $5184$, minimum squared norm $12$,
and value-generated subgroup $4\sqrt2\mathbb Z$ for $\tau|_\Delta$, or
$8\sqrt2\mathbb Z$ for $2\tau|_\Delta$. This is not the actual image:
$12\sqrt2$ is not attained by $\tau$ on $\Delta$.
The preceding proofs retain the triality products and spectral boundary
graph, the flag's degree-$+1$ collapse, and the retained cusp class's
nontrivial three-plane reduction; these exact maps do not supply an
interacting quantum mass-gap result.

For the original analytic fibration $f:X\to\mathbb P^1$, the full
canonical-divisor proof gives
$\omega_X\simeq\mathcal O_X(-2S_2)$ and the graded anticanonical ring
$\bigoplus_{m\ge0}H^0(X,\omega_X^{-m})\cong\mathbb C[U,V]$,
with $\deg U=1$, $\deg V=2$. The complete system $|-2K_X|$
recovers $f$ with the specified projective coordinates. This calculation
uses the original analytic charts and finite fillings, not an
identification of $X$ with $S^6$, and draws no sphere-topology conclusion.

The retained constant $\beta_c=\beta^{\rm part}+c$ now gives a proper
holomorphic deformation over $\operatorname{Im}c<-M$, locally on
relatively compact parameter discs. On a fixed regular fibre set
$T=\operatorname{Im}\tau$, $m=\operatorname{Im}\mu$,
$D=\operatorname{Im}\beta_{c_*}-6m^2/T<0$ and $\delta=c-c_*$.
In the plus-graph convention its exact Beltrami matrix is
\[
 \nu_\delta=\frac{\delta}{2iD+\delta}
                 \begin{pmatrix}0&0\\-m/T&1\end{pmatrix}.
\]
The connecting Kodaira--Spencer class is
$-[\partial_{\zeta_2}\otimes
(d\bar\zeta_2-(m/T)d\bar\zeta_1)/(2iD)]\ne0$.
The intrinsic pencil recovers $f$ also over the dual numbers; its three
fixed critical values prohibit an infinitesimal projective base change.
Consequently this deformation of the whole $X_c$ is nontrivial,
not only its marked regular-fibre deformation. The proof uses no sphere
recognition or CDP conclusion.



The target is a reconstructed manuscript claiming a compact complex
threefold \(X\) and holomorphic map \(f:X\to\mathbb P^1\) with
\[
X\cong_{\mathrm{diff}}S^6,\qquad a(X)=1,\qquad
\langle c_3(T_X),[X]\rangle=2.
\]
As of 6 September 2026, the workbench records the following mathematical
status.  Here \Verified\ means that a complete proof is written in this
workbench; it is not a proxy for a model verdict or for community acceptance.

The 5 September direct-image recheck corrects two written proof steps.
Theorem~\ref{thm:cartier-duality} now constructs its duality morphism
from the analytic trace, not from matching outer terms. The logarithmic
gluing is the real translation
$a_j(s)=\operatorname{Re}\alpha\,v_j+\operatorname{Im}\alpha\,J_sv_j$,
where $J_s=\Pi(s)^{-1}i\Pi(s)$ and $\alpha=\log s_j/(2\pi i)$.
Its full clockwise loop has increment $-v_j$; its mixed sweep integral
is $-E(v_j,\lambda)$. The corrected proof of
Lemma~\ref{lem:picard-finite-norm} and the based calculation in
Proposition~\ref{prop:geometric-pi1} retain their original signs and
coefficients. The bounded recheck verifies these maps, not every
global candidate input. Its before/after sources and primary-source
audit are retained in the direct-image evidence folder.

The period continuation now identifies the exact admissible constants:
for the cusp-regular $\beta=\beta^{\rm part}+c_0$, set
$M=\max_{B\setminus\{p_0\}}D_{\rm part}$, an attained finite maximum.
The marked real period map is everywhere an isomorphism exactly when
$\operatorname{Im}c_0<-M$
(Proposition~\ref{prop:period-exact-admissible-constants}). Omitting the
cusp extension permits $\widetilde\beta=\beta+i\,t\circ\pi$ with
$\widetilde D=D+\operatorname{Re}t$ unbounded above.
The cusp proof uses charts restricted to $Y_\epsilon$, a bound on a
fixed closed subdisc, and the weak inequalities for closed polydiscs
(Lemma~\ref{lem:toric-chart-control}). Its exact opposite-edge maps
retain the three multipliers $e^{-2\pi i\mu_0}$,
$e^{12\pi i\mu_0}$, $e^{-2\pi i(7\mu_0+b_0)}$
with the same-character conventions of
equation~\eqref{eq:cusp-exact-opposite-edge}. These bounded period,
cusp and finite-filling rechecks add no sphere or CDP theorem.

The new higher-wrap comparison is fully written in
Section~\ref{sec:h24-continuation} and in the root draft
\path{s6_higher_rung_24d_preprint.tex}. It gives the octonionic
positive-cone compactification $(B^{25},S^{24})$, the triality flag
$F_4/\operatorname{Spin}(8)$ with integral cohomology ranks $1,2,2,1$,
and the typed peripheral and suspension maps involving
$\mathbb Z/12\hookrightarrow\mathbb Z/24$, $a\mapsto2a$.
The flag facts are attributed to Mare--Willems; the comparison preserves
all coordinate, monodromy, and coefficient choices. The particular
one-copy cusp homotopy coefficient is now zero: the full original
period-corrected derivative frame extends over an explicit collared
$D^2\times T^2$ in one cusp chart times an interval
(Theorem~\ref{h24:thm:cusp-g-zero}). The retained-angle projection of
the same $q_2$ instead gives $H=p_3q_2:X\to S^4$ with
$[H]=\eta_4\eta_5\ne0$ (Theorem~\ref{h24:ret:main}). Its entire
regular inverse image is the original two-torus with the explicitly
calculated normal four-frame and Arf invariant one. Thus $q_2$ is
nonnull and detects a quotient $\pi_6(\mathcal A)/\ker(p_3)_*
\cong\mathbb Z/2$. Smashing $H$ with the specified Hopf map gives
$12\nu_6\ne0$ in $\pi_9(S^6)\cong\mathbb Z/24$, and the
peripheral composite has values $(12,0,12)$ with kernel
$\chi^{-1}(2\mathbb Z/12)$. The full class and exact order of
$q_2$ in its conical target, and its smash powers, are not inferred
from this detected quotient. A completed global higher torus
wrap remains open. The independent arithmetic continuation in
Section~\ref{sec:fifteen-selector-jordan} proves that the fifteen
$p=1201$ projective selector rays span and Jordan-generate the real
three-dimensional symmetric corner, with exact same-residue countercounts.

Section~\ref{esn:sec:shared} adds the explicit even unimodular lattice
$A_5^4\perp D_4\subset N$ of index $72$, with its original $144$ roots,
and the full splitting $O(N)=W(A_5^4D_4)\rtimes L(G)$,
$G\cong GL_2(\mathbb F_3)$
(Theorem~\ref{esn:fa:thm:full-quotient}). This realizes in fixed
coordinate/code markings the abstract Weyl quotient already recorded in
\cite[Table~3]{esn:CDH2014}. The residue six-cycle belongs to the Weyl
group and centralizes this chosen complement
(Theorem~\ref{esn:fa:thm:residue-weyl}); the $D_4$ discriminant action
has kernel $Q_8$ in $G$ and image the full triality $S_3$.
The fixed-shell divisor bijection and finite-sample Fourier map retain
$p,R,P$: $A_3(2521)=0$, $A_3(3361)=6$, with unordered counts $0,3$,
despite equal residues modulo $840$.

The later quaternionic extension is treated separately in
Section~\ref{sec:q444-reconciliation}: the two Lorentz thirteen-slices
coincide as embedded matrix spaces, but the supplied unconstrained scalar
model can leave their trace-free slice. The explicit witness gives
$(u^*b)^{\cdot\cdot}=-2\kappa$ at $u^*b=(u^*b)^{\cdot}=0$.
There are also global interacting invariant families within the original
equations: fixing $W\subset u^\perp$, retaining $b\in W$ and
$\Phi W\subset W$, with tangent velocities, gives configuration dimensions
six and eight and identically zero reaction. A separate fourth-order
escape example has $\rho^{(j)}(0)=0$ for $0\leq j\leq3$ but
$\rho^{(4)}(0)=-12\kappa$. These results distinguish exact invariant
families from instantaneous tangency tests. The families do not exhaust
the invariant motions: Theorem~\ref{q444:moving-family} constructs
periodic moving-kernel solutions outside every fixed-$W$ family.
Their conserved nonzero colour moment gives an explicit Gauss residual
for the zero-connection gauge lift, not an obstruction to every possible
gauge-field lift. The complete scalar vacuum diagram and full metric
Hessian are also proved in Theorems~\ref{q444:vacua} and
\ref{q444:vacuum-hessian}; their gauge-orbit and quantum counterparts are
not inferred from this calculation.
The full equations and the different constrained equations are both given;
the original candidate's complex-structure claims do not follow from this
quaternionic cone calculation.

The full positive fifteen-space also has a split positive observation to
$\operatorname{Herm}_2(\mathbb C)$ with exact image condition
$|\rho|^2\leq\det G$ (Proposition~\ref{q444:full-observation}).
Every permitted $(G,\rho)$ occurs; null $G$ forces $\rho=0$, while
$G=\epsilon I_2$, $\rho=\epsilon/2$ permits positive lifts for every
$\epsilon>0$. For trace-free scalar matrices the distinct spectral-floor
map \eqref{q444:spectral-floor} gives this positive observation with the
original matrix $C=PP^*$ retained. These exact maps do not impose a
positive uniform spectral threshold or identify the coordinates as
gauge-invariant physical momenta.

The complete retained-period metric fibres are now proved in
Theorem~\ref{q444:metric-fibre}: they have dimensions twelve over $H_3$
and seventeen over $G$. Every smooth positive observed trajectory has
an equivariant lift with the original periods unchanged. An explicit
$G=\operatorname{diag}(\alpha_{\rm obs},\beta_{\rm obs}y^{-p})$
approaches a nonzero null block; its relative metric condition number
is asymptotic to
$2\max\{2\alpha_{\rm obs}/3,3\lambda\}y^{p+2}/\beta_{\rm obs}$.
Propositions~\ref{q444:metric-condition}--\ref{q444:rank-loss-metric}
give the exact inverse and bounds. Positivity and period monodromy
alone consequently do not select an observed trajectory.

Theorem~\ref{q444:temporal-no-lift} strengthens the rotating-orbit
obstruction to every homogeneous spatial connection in temporal gauge.
Proposition~\ref{q444:gauss-inertia} retains the full arbitrary-temporal
Gauss equation without inferring the remaining equations from it.
Theorem~\ref{q444:radial-gauge} constructs a different interacting
two-variable family satisfying every field equation.
Finally, Proposition~\ref{q444:vacuum-average} identifies
$\int\det(PgAg^*P^*)\,d\mu
=V_\mu\operatorname{tr}\operatorname{adj}(A)=V_\mu\|b\|^2$
on the positive rank-one boundary and exhibits gauge-inequivalent
scalar minima on which $\|b\|^2$ ranges continuously.
The fixed-mark determinant itself varies within a gauge orbit.

Theorems~\ref{q444:real-no-lift} and \ref{q444:static-complex-no-lift}
now exclude all homogeneous real connections with arbitrary time
dependence, and all complex connections constant in the spatial
constant-matter coordinates, respectively, for the literal rotating orbit.
Theorem~\ref{q444:cx:thm:block} additionally excludes the genuinely
complex rotating-pair class $C,Z_1\in\mathfrak p$, $Z_2\in\mathfrak k$
and its constant $SO(2)$ spatial saturation, where the block splitting
is defined by the fixed matter line $\mathbb Cq$. No unrestricted
pair is assumed to admit that splitting. For arbitrary spatial colour
matrices in the rotating-pair ansatz, Lemma~\ref{q444:complex-gauss-inverse}
proves $g^2\mathcal I+2\mathcal K>0$ whenever $[Z_1,Z_2]\ne0$, a
condition forced by every solution. Theorem~\ref{q444:complex-exact-zero-set}
gives its exact Gauss/frequency inverse and all four residuals
\eqref{q444:complex-full-residuals}. Whether their zero set is empty,
and the existence of still more general time-dependent complex lifts,
remain open; neither is decided by a scalar contraction alone.
The complete scalar equivalence in
Proposition~\ref{q444:complex-full-scalar-equivalence} and the
strict stabilizer signature constraint of
Theorem~\ref{q444:complex-stabilizer-spacelike} hold in every
complex coordinate, not only in the excluded block class.
Proposition~\ref{q444:complex-stress-obstruction} adds the exact
two-by-two spatial stress equation and its sum-of-squares discriminant.

Theorem~\ref{q444:finite-coupling-fibre} now computes the complete
coupling fibre over every unrestricted ordered spatial pair as the
admissible zeros of a nonnegative degree-twenty-eight polynomial.
There are at most fourteen distinct positive $g^2$ values, and its
explicit inverse retains both signs of $\omega$ and every temporal
coordinate. Proposition~\ref{q444:finite-coupling-domain-test} gives
three exact vector tests for an identically empty denominator domain.
Theorem~\ref{q444:generic-projection-obstruction} additionally places
the full spatial projection inside the zero set of a specific nonzero
Sylvester determinant polynomial on all sixteen coordinates. This
set is closed, nowhere dense and measure zero; an explicit full pair
certifies the determinant is nonzero. Proposition~\ref{q444:twocross:prop:exceptional-inclusion} places the
entire eight-dimensional two-cross spatial subspace inside that zero set,
and Theorem~\ref{q444:twocross:thm:exclusion} excludes all full solutions
there. Proposition~\ref{q444:twocross:quantitative} gives a polynomial
residual certificate and a positive uniform lower bound for spatial error
on its bounded scalar--Gauss zero locus at coupling bounded away from zero.
The full scalar--Gauss zero locus is now completely parameterized by
Theorems~\ref{sgd:parametrization} and~\ref{sgd:literal-addon:quotient}:
\[
 \mathcal D=\{(u,v_0,y):u>0,\ d=uv_0-\delta^2-y^2>0,
                       \ c=4v_0-5u>0\},
\]
\[
 \begin{gathered}
 \mathcal S_{\rm sg}\cong
 \bigl(\mathcal D\times\{1,-1\}^2\times U(1)^2\times U(1)\bigr)
 /\langle\iota\rangle,\\
 \iota:(p_+,p_-,\zeta)\mapsto(-p_+,-p_-,-\zeta),
 \qquad \delta=\sqrt{2/3}.
 \end{gathered}
\]
where $\iota$ fixes the other coordinates and $t=72d/c$.
The smooth refinement in Corollary~\ref{sgd:smooth-addon:levels} gives
\[
 \mathcal S_{\rm sg}\cong\{1,-1\}^2\times(0,\infty)
                  \times\mathbb R^2\times\mathbb T^3.
\]
These are four embedded six-dimensional literal components. The original
coupling is the displayed positive coordinate and is a surjective
submersion; every fixed $t=\tau>0$ has exactly four embedded
five-dimensional components $\mathbb R^2\times\mathbb T^3$.
For every prescribed original $\tau=g^2>0$,
\[
 \rho=\frac{\tau+\sqrt{\tau^2+4608\delta^2}}{96\delta},\qquad
 M_\rho=\delta\begin{pmatrix}\rho&1\\1&2\rho\end{pmatrix}
\]
gives $t=\tau$ by Corollary~\ref{sgd:every-coupling}, with every
scalar and Gauss entry zero and both original frequency signs retained.
The two spatial residuals cannot vanish simultaneously in this class.
Existence on the remaining exceptional locus is \Open. None of these finite-ansatz results
is a theorem about arbitrary time-dependent gauge connections or
the interacting continuum quantum spectrum.

The different sector in Proposition~\ref{q444:vk-sector} supplies
exact full-gauge solutions. Its periodic family
(Theorem~\ref{q444:vk-periodic}) has $\omega\to0$ and
$E-V_0=9h_0^2\omega^2+O(\omega^4)$ at fixed spatial volume;
its amplitude tends to $3h_0$, not zero.
Proposition~\ref{q444:vk-wave} proves nonlinear null waves and their
finite-volume energy. The entire transverse gauge spectrum is given in
Theorem~\ref{q444:gauge-spectrum}. The gauge-invariant conserved vector
current in Proposition~\ref{q444:current-mass} detects a branch with
$\mu_{\rm d}^2=(2g^2/3)\|b\|^2$, exactly proportional to the preceding
colour average. Its linear current coefficient vanishes at $b=0$.
These are proved classical equations and observables, not quantum
spectral conclusions.

There is now a separate quantum result for the complete homogeneous
gauge--scalar mechanics, not inferred from those classical families.
Section~\ref{qh:chapter} retains all thirty-eight real configuration
coordinates. The exact form bound \eqref{qh:full-form-bound}, with
its original $\hbar,g^2,\mathcal V_3,\Gamma,m^2,\kappa$, proves compact
resolvent. Theorems~\ref{qh:hq-physical} and \ref{qh:hq-positive}
impose the spatially constant Gauss constraints and prove a unique
invariant ground state; Corollary~\ref{qh:hq-ground} gives
$E^{\rm phys}_1-E_0>0$ for fixed positive volume and couplings.
This is an exact finite-dimensional interacting quantum result,
not a continuum Yang--Mills gap or a uniform infinite-volume estimate.
The classical embedding and the failure of the automatic product
quantum-subsector map are proved explicitly in
Proposition~\ref{qh:hq-embedding} and the following Fourier calculation.
The bosonic commutator-matrix discreteness mechanism is attributed
to Simon's classical work, not claimed as new.
Section~\ref{fq:section} constructs a further exact correspondence for
the retained two-Fourier-coordinate mechanics. Its true positive
quartic ground state $\Omega_q$ gives $Jf=f\Omega_q$, zero Berry
connection and the nonzero Born--Huang term in \eqref{fq:effective}.
The full residual \eqref{fq:intertwining} prevents this moving
subspace from being reducing, but Theorem~\ref{fq:dynamics} bounds
the difference of its physical-time evolutions by the explicit
$|t|\mathcal E(E)$ in \eqref{fq:time-error}. The sufficient regime
\eqref{fq:regime} is proved by min--max and a variational bound,
not an assumed spectral gap. No Gaussian is substituted for the
quartic ground state. The error is of order $\mathcal V_3^{-1/2}$
at fixed $\hbar,\Gamma,L_1,m^2,E,t$; it does not control a simultaneous
limit in which $\omega_*\to0$, nor the full continuum theory.

Theorem~\ref{q444:hermitian-cometric} identifies the original weight
$W$ as a lattice cometric. Complex-structure compatibility removes
exactly its twelve anti-complex-structure fibre directions and gives
$h_3=2H_3^{-1}$. The quotient in
Theorem~\ref{q444:metric-submersion} has $h_2=2G^{-1}$ and intertwines
the Laplacians on the exact pulled-back character subspace.
Proposition~\ref{q444:native-period-spectrum} retains the native
Euclidean metric and all inherited period corrections: $H_3=2I_3$,
$G=2PP^*$, $\lambda_1^{(3)}\sim4\pi^2/y^2$ and
$\lambda_1^{(2)}\sim8\pi^2/y^2$.
The full marked spectrum recovers the lattice cometric even though
the fixed observed matrix does not recover the changing periods.

\begin{enumerate}[leftmargin=1.8em]
\item \Verified\ Let
\(\Delta\simeq(\mathbb Z/3)*(\mathbb Z/4)\) act on
\(\mathfrak h_z\), let \(g_1,g_2\) be its elliptic generators with chosen
fixed points \(z_1,z_2\), and put \(g_0=(g_1g_2)^{-1}\).  For the normalized
period lift \(\tau:\mathfrak h_z\to\mathfrak h\), define
\[
 \varphi:\Delta\longrightarrow\operatorname{PSL}_2(\mathbb Z),
 \qquad \tau\circ g=\varphi(g)\circ\tau.
\]
To avoid confusing the modular transformations with the finite reduced
fibres \(S_j\) below, write
\[
 \mathfrak s_1(w)=\frac{w-1}{w},\qquad
 \mathfrak s_2(w)=-\frac1w,\qquad
 \mathfrak p(w)=w-1.
\]
The finite generators are normalized simultaneously by
\[
 g_1'(z_1)=e^{-2\pi i/3},\qquad g_2'(z_2)=-i,\qquad
 g_0\ \text{parabolic},
\]
and the exact modular deck certificate is
\[
 \varphi(g_1)=\mathfrak s_1,\qquad
 \varphi(g_2)=\mathfrak s_2,\qquad
 \varphi(g_0)=\mathfrak p,\qquad
 \mathfrak s_1\mathfrak s_2\mathfrak p=1.
\]
The sign at \(z_2\) is forced by the simultaneous parabolic normalization;
the quadratic ramification calculation alone gives only
\((g_2'(z_2))^2=-1\).  The cusp normalization proof now records its
covering levels explicitly: it chooses a \(g_0\)-stable component
\(U_0\subset\pi^{-1}(D_0^*)\) and a component
\(K\subset j^{-1}(\{|w|>R\})\) containing \(\tau(U_0)\).  For
\(P_{\rm cusp}:=\varphi(g_0)\), the auxiliary conjugation
\[
 \tau'=\gamma_0^{-1}\tau|_{U_0},\qquad
 P'=\gamma_0^{-1}P_{\rm cusp}\gamma_0,\qquad
 \gamma_0(\infty)=c,
\]
where \(c\) is the rational boundary cusp of \(K\), proves that
\(P'\ne1\) and preserves the selected cusp at infinity, hence
is parabolic.  The proof then discards \(\gamma_0\), returns to the original
\(\mathfrak s_1\)-normalized frame, performs only the allowed final
\(\mathfrak s_1\)-power conjugation, and obtains the displayed modular
certificate.  Only after that return does it define anew the cusp
exponential
\[
 q_{\rm cusp}:=e^{2\pi i\tau}=t_cu_1(t_c),\qquad u_1(0)\ne0;
\]
this \(q_{\rm cusp}\) is distinct from the integral class \(q\) below.
The resulting nondegenerate period family descends, and the middle quotient
\[
 J=(\mathcal T|_{\mathfrak h_z^\circ})/\Delta
 \longrightarrow B_{\rm mid}:=\mathfrak h_z^\circ/\Delta
\]
is a proper holomorphic submersion with compact torus fibres.  The toric
cusp filling, the two finite logarithmic fillings, all three
punctured-overlap biholomorphisms, and the global gluing produce a connected
compact Hausdorff smooth complex threefold and a proper surjective
holomorphic map with connected fibres
\[
 f:X\longrightarrow B\simeq\mathbb P^1.
\]
Its full scheme-theoretic cusp fibre \(X_{p_0}=W\) is reduced, irreducible,
nonnormal, and normal crossings, with local equations
\[
 z_0,\qquad z_0z_1,\qquad z_0z_1z_2,
\]
and normalization \(\nu:E_0\simeq\mathrm{dP}_6\to W\); it is not a
simple-normal-crossings divisor.  For the finite fibres, let
\[
 A_0^{(j)}=(\Lambda\otimes_{\mathbb Z}\mathbb R)/\Lambda\simeq F
\]
be the marked central torus on the covering family, and put
\[
 v_1=\widehat\gamma+2\widehat u-4\widehat w,\qquad
 v_2=-(\widehat\gamma+3\widehat u-3\widehat w),
\]
so that
\[
 S_j=A_0^{(j)}/\langle x\mapsto A_jx+v_j/m_j\rangle,\qquad
 f^*(p_j)=m_jS_j,\qquad (m_1,m_2)=(3,4).
\]
The quotient covering and its codomain are, explicitly,
\[
 \pi_j:A_0^{(j)}\longrightarrow S_j,\qquad x\longmapsto[x].
\]
The radial retraction \(N_j\to S_j\) identifies finite-fibre specialization
exactly with
\[
 \operatorname{sp}_{p_j}^r=\pi_j^*:
 H^r(S_j;\mathbb Z)\longrightarrow H^r(A_0^{(j)};\mathbb Z)
 \xrightarrow[\text{marking}]{\ \sim\ }H^r(F;\mathbb Z),
\]
with no unrecorded conjugacy or index.  This stalk identification is
\cref{lem:finite-filling-stalk-map}; see also
\cref{lem:triangle-generator-orientations,prop:period-input-audit,%
thm:candidate-cusp-geometry,prop:candidate-cusp-overlap,%
lem:middle-period-proper,prop:candidate-finite-fillings,%
thm:candidate-global-geometry}.

\item \Verified\ The monodromy, cusp invariant lattices, global
coinvariants, finite-filling stalk maps, terminal group presentation, and
terminal degree-two lattice are proved over \(\mathbb Z\) in
\cref{prop:monodromy-certificate,thm:cusp-lattices,thm:coinvariants,%
prop:terminal-group,prop:terminal-degree-two,prop:terminal-all-degree,prop:finite-indices}.

\item \Verified\ The intrinsic integral nearby-cycle specialization at the
cusp is injective with saturated image
\[
 \operatorname{im}(\operatorname{sp}^q)
   =H^q(F;\mathbb Z)^{T_0}\qquad(0\le q\le4),
\]
where the degree-two invariant class is
$q:=uw+6\gamma\delta\in\bigwedge^2V=H^2(F;\mathbb Z)$,
and the finite stalks and parabolic sheaf calculation give the complete
global Leray page.  The canonical truncation triangle identifies the
explicit Cech discrepancy cocycle with Leray \(d_2\), and multiplicativity
then gives
\[
 d_2(12\gamma)=\omega,\qquad
 d_2(2q)=\omega[\delta],\qquad
 d_2(2\gamma uw)=\omega[\gamma\delta].
\]
Consequently
\[
 H^1(X;\mathbb Z)=H^2(X;\mathbb Z)=H^3(X;\mathbb Z)=0,
 \qquad H_2(X;\mathbb Z)=0;
\]
see
\cref{thm:p0-integral-specialization,prop:leray-e2,lem:first-transgression,lem:higher-transgressions,thm:h2-vanishing}.

\item \Verified\ The candidate's independent toric-collapse construction is
continuous across every stratum and descends equivariantly to a strong
deformation retraction $r:N'_0\to W$.  Its fixed-fibre restriction
$c=r|_F$ has the exact primitive integral sequence
\[
0\longrightarrow\operatorname{im}(\textstyle\bigwedge^qM_0-I)
\longrightarrow\textstyle\bigwedge^q\Lambda
\xrightarrow{c_*}H_q(W;\mathbb Z)\longrightarrow0,
\]
and proper-base-change functoriality gives, in every degree,
\[
 c^*=i_F^*r^*=\operatorname{sp}^q.
\]
Thus the former C1--C3 obligations are closed; see
\cref{prop:verified-cellular-collapse,prop:verified-collapse-homology,%
thm:collapse-equals-specialization}.

\item \Verified\ With zero-section whiskers and the fixed clockwise
meridians, the actual overlap inclusions give
\[
 \pi_1(X)=
 \langle c,x,y\mid c\text{ central},\ xy=1,\ x^3=c,\ y^4=c^{-1}\rangle
 \simeq1.
\]
Together with the preceding integral cohomology this proves
\[
 X\cong_{\mathrm{diff}}S^6,\qquad
 \langle c_3(T_X),[X]\rangle=2;
\]
see \cref{prop:geometric-pi1,cor:candidate-sphere}.

\item \Verified\ The very-general fibre has
\[
 \eta:=q=u\wedge w+6\gamma\wedge\delta,\qquad
 \operatorname{NS}(F_b)=\mathbb Z\eta,\qquad a(F_b)=0.
\]
The square is the oriented wedge identity
\[
 \eta\wedge\eta
 =12\,\gamma\wedge u\wedge w\wedge\delta,
\]
so, for the orientation determined by the ordered marking
\((\gamma,u,w,\delta)\) and for the complex orientation respectively,
\[
 \int_{(F_b,\operatorname{or}_{\rm mark})}\eta\wedge\eta=12,
 \qquad
 \int_{(F_b,\operatorname{or}_{\mathbb C})}\eta\wedge\eta=-12.
\]
The repaired meromorphic-descent argument gives
\[
 \mathcal M(X)=f^*\mathbb C(t),\qquad a(X)=1,
\]
with \(f\) the algebraic reduction; see
\cref{prop:very-general-torus-audit,thm:algebraic-dimension-audit}.

\item \Refuted\ The p.~692 sheaf-theoretic inference of
Campana--Demailly--Peternell is false in the stated generality.  If
\(
\mathscr K_S=\ker(\Omega_X^1|_S\to\widetilde\Omega_S^1)
\), then
\[
0\longrightarrow N^*_{S/X}\longrightarrow\mathscr K_S
\longrightarrow\operatorname{tors}\Omega_S^1\longrightarrow0,
\]
and the two printed vanishings do not control sections mapping to the
conductor-supported quotient.  The proper nodal-cubic-times-
\(\mathbb P^1\) family in \cref{prop:nodal-counterexample} is an explicit
counterexample to that inference.

\item \Verified\ The conductor section, nonnormal Cartier-divisor duality,
flatness, and top-degree analytic base change are complete in
\cref{prop:conductor,thm:cartier-duality,prop:cartier-duality-absolute,%
lem:curve-flatness,prop:analytic-base-change,cor:surface-top-base-change}.
On the cusp model the invariant nowhere-vanishing form
\[
 \Omega=\frac{dx_1\wedge dx_2\wedge dt}{x_1x_2}
\]
first proves intrinsically that
\(\omega_{N_0}\simeq\mathcal O_{N_0}\).  After the global gluing, the open
embedding \(N_0\hookrightarrow X\) gives the exact restriction chain
\[
 \omega_X|_{N_0}\simeq\omega_{N_0}\simeq\mathcal O_{N_0}.
\]
Let \(\nu:E_0\simeq\mathrm{dP}_6\to W\) be the normalization.
Consequently, for every \(L\in\operatorname{Pic}^0(X)\),
\[
 \nu^*((L^\vee\otimes\omega_X)|_W)
 \simeq\mathcal O_{\mathrm{dP}_6},
\]
and the conductor--duality--base-change chain gives
\[
 R^2f_*(T_X\otimes L)\ne0
\]
without using global topology; see
\cref{thm:candidate-cusp-geometry,thm:candidate-global-geometry,%
prop:pic0-pullback,cor:candidate-pic0-route}.

\item \Verified\ For arbitrary line bundles the exact trivial-pullback
obstruction is
\[
 r=\nu^*i^*:\operatorname{Pic}(X)\longrightarrow
 \operatorname{Pic}(\mathrm{dP}_6).
\]
The conductor argument applies uniformly to all line bundles exactly when
$r=0$; see \cref{prop:all-pic-restriction}.  For the constructed candidate,
the Leray edge image of every $c_1(L)$ lies in
$E_\infty^{0,2}\subseteq\ker d_2^{0,2}$, while
$d_2^{0,2}(2q)=\omega[\delta]$ is an isomorphism.  Proper-base-change
naturality therefore gives $c_1(L|_W)=0$, and injectivity of $c_1$ on
$\operatorname{Pic}(\mathrm{dP}_6)$ forces $r=0$; see
\cref{prop:leray-edge-restriction}.  Combining this with the analytic inputs
gives
\[
 R^2f_*(T_X\otimes L)\ne0
 \qquad\text{for every }L\in\operatorname{Pic}(X).
\]
Thus the topology-free Pic$^0$ route and the all-Pic route are compatible
rather than competing statuses.  There are now three all-Pic certificates.
The first is the preceding Leray-edge proof.  The second uses no global
Leray differential: the order-four specialization lattice writes
$c_1(L|_F)=2nq$, and the pair-of-pants crossed defect has boundary values
\[
 \bar d_1=4n[\delta],\qquad
 \bar d_2=-3n[\delta],\qquad
 \bar d_0=0.
\]
The exact boundary relation forces $n=0$, and local integral cusp
specialization then forces $r(L)=0$; see
\cref{lem:picard-crossed-defect,lem:picard-finite-norm,%
lem:picard-cusp-defect,prop:direct-boundary-picard}.  This second proof is
independent of the global transgressions, but it still uses one local cusp
bounding chain and the local specialization theorem.

\item \Verified\ The local cusp Picard map itself is nonzero.  The explicit
equivariant Cartier divisor
\[
 \mathcal D_Q=\sum_{v\in\mathbb Z^2}(v_1^2-6v_2^2)E_v
\]
descends to $\mathcal L_Q\in\operatorname{Pic}(N_0)$ with
\[
 c_1(\mathcal L_Q|_{F_t})=-2q,\qquad
 \bigl(\deg_{C_i}\nu^*(\mathcal L_Q|_W)\bigr)_{i=1}^6
 =(-12,2,0,-12,2,0).
\]
Thus neither the cusp geometry nor analytic relative-Picard
representability alone can force normalized central restriction to vanish;
see \cref{prop:local-cusp-picard-obstruction}.

\item \Verified\ After the global coinvariant calculation writes
\[
 E:=c_1(L|_F)=kq\in\operatorname{Alt}^2(\Lambda,\mathbb Z)
       =\bigwedge^2V,
\]
the third all-Pic certificate completes the argument by holomorphic factors
and toric support functions.  Put
\(V_{\mathbb R}=\operatorname{Hom}(\Lambda,\mathbb R)\), let
\(P\subset B^\circ\) be the compact pair-of-pants core with clockwise
boundary meridians \(\rho_1,\rho_2,\rho_0\), and let
\(d:\pi_1(P)\to V_{\mathbb R}\) be the factor-of-automorphy crossed
homomorphism, put \(d_i=d(\rho_i)\), and let a bar denote passage to
\(V_G\otimes\mathbb R=\mathbb R[\delta]\).  Logarithms of the actual
\(\mathcal O^*\)-factors give
\[
 \bar d_1+\bar d_2+\bar d_0=0,\qquad
 \bar d_1=2k[\delta],\quad
 \bar d_2=-\tfrac32k[\delta],\quad
 \bar d_0=0.
\]
The finite norm identities now include the fully typed factor-extension
repair (5.17-A14c)--(5.17-A15), rather than identifying maps living on
different covering levels.  Put
\[
 \widetilde J_j=\Delta_j\times\mathbb C^2,
 \qquad r_j:\widetilde J_j\longrightarrow J'_j,
\]
where \(r_j\) is the lattice quotient with deck translations
\(\tau_\lambda\).  On this vector cover, and only there, define
\[
 Q_{j,\log}=\widetilde g_j^{\,v_j/m_j},\qquad
 Q_{j,\mathrm{lin}}=\widetilde g_j^{\,0}.
\]
Their exact relations are
\[
 Q_{j,\log}\tau_\lambda Q_{j,\log}^{-1}=\tau_{A_j\lambda},
 \quad (Q_{j,\log})^{m_j}=\tau_{v_j},
 \qquad
 Q_{j,\mathrm{lin}}\tau_\lambda Q_{j,\mathrm{lin}}^{-1}
 =\tau_{A_j\lambda},
 \quad (Q_{j,\mathrm{lin}})^{m_j}=1.
\]
Let \(q_{j,\log},q_{j,\mathrm{lin}}\) be their descended order-\(m_j\)
automorphisms of \(J'_j\), put
\(J_j^{\prime *}=J'_j|_{\Delta_j^*}\), and define
\[
 p_{j,\log}:J_j^{\prime *}\to
 N_j^*:=J_j^{\prime *}/\langle q_{j,\log}\rangle,
 \qquad
 p_{j,\mathrm{lin}}:J_j^{\prime *}\to
 J^*:=J_j^{\prime *}/\langle q_{j,\mathrm{lin}}\rangle.
\]
The branch-independent fibre translation
\(h_j:J_j^{\prime *}\to J_j^{\prime *}\) and
its quotient biholomorphism \(\bar h_j:N_j^*\to J^*\) satisfy
\[
 h_jq_{j,\log}h_j^{-1}=q_{j,\mathrm{lin}},
 \qquad p_{j,\mathrm{lin}}\circ h_j=\bar h_j\circ p_{j,\log},
\]
and the gluing gives the typed line-bundle comparison
\[
 p_{j,\log}^*(L|_{N_j^*})
 \simeq h_j^*p_{j,\mathrm{lin}}^*(L|_{J^*}).
\]

For the filled factor extension return to the vector cover.  With
\[
 P_{j,\log}:J'_j\to N_j,\qquad
 \pi_j^{\rm vec}:=P_{j,\log}\circ r_j:\widetilde J_j\to N_j,
 \qquad
 \Gamma_j^{\log}
 =\langle\tau_\lambda,Q_{j,\log}:\lambda\in\Lambda\rangle,
\]
where \(P_{j,\log}\) is the full quotient and
\(P_{j,\log}|_{J_j^{\prime *}}=p_{j,\log}\), pull back
\(L|_{N_j}\) by \(\pi_j^{\rm vec}\), push its integral factor
extension out along \(\mathbb Z\hookrightarrow\mathbb R\), and choose
normalized vertical lifts \(U_\lambda^{(j)}\).  For
\(\alpha\in\Gamma_j^{\log}\), define
\(A_\alpha\in\operatorname{Aut}(\Lambda)\) and
\(T_\alpha\in\operatorname{Aut}(V_{\mathbb R})\) by
\[
 \alpha\tau_\lambda\alpha^{-1}=\tau_{A_\alpha\lambda},
 \qquad T_\alpha=(A_\alpha^{-1})^{\mathsf t}.
\]
For a factor-extension lift \(\widetilde G_\alpha\), define
\(\widetilde d_j(\alpha)\in V_{\mathbb R}\) by
\[
 \widetilde G_\alpha U_{A_\alpha^{-1}\lambda}^{(j)}
 \widetilde G_\alpha^{-1}
 =\mathbf c\!\left(\langle\widetilde d_j(\alpha),\lambda\rangle\right)
 U_\lambda^{(j)}.
\]
The normalized product gives
\[
 \widetilde d_j(\alpha\beta)
 =\widetilde d_j(\alpha)+T_\alpha\widetilde d_j(\beta).
\]
Compatible basepoints in the quotient square identify the clockwise
boundary lift successively with \(q_{j,\mathrm{lin}}\), \(q_{j,\log}\), and
finally the vector-cover deck transformation \(Q_{j,\log}\).  The induced
extension isomorphism then proves, for some \(a_j\in V_{\mathbb R}\),
\[
 \widetilde d_j(Q_{j,\log})=d_j+(T_j-I)a_j,
 \qquad
 \widetilde d_j(\tau_v)=\iota_vE,\qquad
 (\iota_vE)(\lambda)=E(v,\lambda).
\]
Consequently, with \(N_{T_j}=I+T_j+\cdots+T_j^{m_j-1}\),
\[
 N_{T_j}d_j
 =N_{T_j}\widetilde d_j(Q_{j,\log})
 =\widetilde d_j((Q_{j,\log})^{m_j})
 =\widetilde d_j(\tau_{v_j})=\iota_{v_j}E,
\]
because \(N_{T_j}(T_j-I)=0\).  This proves the finite boundary values,
including their signs, inside the actual filled logarithmic deck group and
closes the factor-extension gap; the certificate therefore remains
{\Verified}.  Changes of line-bundle frame, central lift, normalized
vertical system, or logarithm branch alter only the displayed principal
character term and leave \(h_j,\bar h_j\), the quotient square, and
(5.17-A15) unchanged.
The cusp equality follows from the argument principle in a Stein chart, not
from a Chern evaluation on a bounding chain; hence \(k=0\).  For adjacent
rays in a maximal cone
\(\sigma=\{v,w,u\}\), the actual toric-open inclusions are
\[
 U_{\rho_v}=D(z_wz_u)\subset U_\sigma,\qquad
 U_{\rho_w}=D(z_vz_u)\subset U_\sigma,\qquad
 U_{\rho_v}\cap U_{\rho_w}=U_0.
\]
The contractible Stein comparison chart
\[
 S_\sigma=\{(z_v,z_w,z_u)\in\mathbb C^3:
               |z_vz_wz_u|<\epsilon\}
\]
has \(\operatorname{Pic}(S_\sigma)=0\), so its frames give zero horizontal
winding along the connected one-skeleton.  For the support-function Cartier
divisor \(D_a\), the exact deck-translate identity
\[
 \Psi_\lambda^*D_a-D_a=\operatorname{div}(\chi^{m_\lambda})
\]
shows that multiplication by \(\chi^{m_\lambda}\) extends from the dense
torus to an isomorphism
\[
 \mu^+_\lambda:\Psi_\lambda^*\mathcal O(D_a)
 \xrightarrow{\ \sim\ }\mathcal O(D_a).
\]
The map on the negative divisor required by the descent tensor product is
the inverse-dual
\[
 \mu^-_\lambda=((\mu^+_\lambda)^{-1})^\vee:
 \Psi_\lambda^*\mathcal O(-D_a)\longrightarrow\mathcal O(-D_a),\qquad
 s\longmapsto\chi^{-m_\lambda}s.
\]
Here the cusp shear is the typed lattice isomorphism
\[
 \Lambda_{\rm tor}:=\ker(M_0-I)=\operatorname{im}(M_0-I)
   =\mathbb Z\langle\widehat w,\widehat\delta\rangle,
 \qquad \overline\Lambda:=\Lambda/\Lambda_{\rm tor},
\]
\[
 B_0:\overline\Lambda\xrightarrow{\ \sim\ }\Lambda_{\rm tor},
 \qquad
 B_0\overline{\widehat\gamma}=-\widehat\delta,
 \qquad B_0\overline{\widehat u}=\widehat w.
\]
Thus its matrix is
\(\left(\begin{smallmatrix}0&1\\-1&0\end{smallmatrix}\right)\)
from the ordered basis
\((\overline{\widehat\gamma},\overline{\widehat u})\) to
\((\widehat w,\widehat\delta)\); in particular \(B_0\) is a map between
two named lattices, not an untyped endomorphism, and \(\det B_0=1\).
For \(M=L|_{N_0}\), if \((d_1,d_2,d_3)\) are the independent degrees of
\(\nu^*(M|_W)\), then
\[
 \bigl(\deg_{C_i}\nu^*(M|_W)\bigr)_{i=1}^6
 =(d_1,d_2,d_3,d_1,d_2,d_3),
\]
and their support function has polarization
\[
 B\in\operatorname{Sym}^2(\Lambda_{\rm tor}^*),\qquad
 [B]_{(\widehat w,\widehat\delta)}
 =\begin{pmatrix}d_2+d_3&d_3\\d_3&d_1+d_3\end{pmatrix}.
\]
For the Appell--Humbert form
\(E_M:=c_1(M|_F)\in\operatorname{Alt}^2(\Lambda,\mathbb Z)\), one has
\[
 E_M((\lambda,0),(0,m))=-B(B_0\lambda,m)
 \quad
 (\lambda\in\overline\Lambda,\ m\in\Lambda_{\rm tor}).
\]
The already proved antecedent \(k=0\) gives
\(E_M=c_1(L|_F)=0\); only then does the unimodularity of \(B_0\)
force \(B=0\), then \(d_1=d_2=d_3=0\), and hence
\(\nu^*i^*L\simeq\mathcal O_{\mathrm{dP}_6}\).  Beyond the preceding
coinvariant input, this proof uses neither a global integral Leray
transgression, local integral nearby-cycle specialization, nor Chern-class
evaluation on a cusp bounding chain; see
\cref{prop:analytic-ah-boundary,lem:analytic-support-specialization,%
thm:analytic-all-picard}.

\item \Refuted\ The separate split-extension calculation imported from
CDP98 replaces monodromy coinvariants by the entire fibre homology.  The
correct formula is
\[
 (K\rtimes Q)^{\mathrm{ab}}
 \simeq Q^{\mathrm{ab}}\oplus(K^{\mathrm{ab}})_Q.
\]

\item \Verified\ A corrected replacement theorem now survives both printed
defects.  If one \(L_0\in\operatorname{Pic}(X)\) satisfies
\[
 H^2(X,T_X\otimes L_0)=0,
\]
then generic endpoint vanishing, semicontinuity, and
Hirzebruch--Riemann--Roch give
\[
 H^i(X,T_X\otimes M)=0\ (i\ne1),\qquad
 \chi(X,T_X\otimes M)=\tfrac12c_3(X)\le0
\]
for general \(M\).  It is enough that
\(R^2f_*(T_X\otimes L)=0\) for one \(L\).  For reduced full fibres, the
invalid p.~692 shortcut is replaced exactly by injectivity of
\[
 \delta_c:H^0(W_c,\operatorname{tors}\Omega^1_{W_c}\otimes A_c)
 \longrightarrow H^1(W_c,N^*_{W_c/X}\otimes A_c)
\]
together with the two endpoint vanishings.  See
\cref{thm:corrected-cdp-global,cor:corrected-cdp-relative,%
 thm:corrected-cdp-fibrewise}.  The constructed candidate fails the global and
 relative conditions for separate reasons.  Since \(c_3(X)=2\), the
 contrapositive of \cref{thm:corrected-cdp-global} gives
 \(H^2(X,T_X\otimes L_0)\ne0\) for every \(L_0\).  Independently, its
 conductor makes \(R^2f_*(T_X\otimes L)\ne0\) for every \(L\), so the stronger
 relative sufficient condition also fails.  The original CDP20 theorem
 therefore remains refuted.

\item \Verified\ The constructed \(X\) satisfies the stated hypotheses of
CDP20 Theorem~2.2 but has both \(c_3(X)=2>0\) and \(a(X)=1\).  Thus it is a
counterexample to that theorem and its ensuing \(S^6\) corollary as stated.
This theorem-level conclusion uses the complete construction, topology, and
algebraic-reduction proofs above; it is not inferred merely from the two
printed proof-route defects.

\item \Verified\ The candidate's Proposition~10.8 conclusion holds.  The
workbench proves
\[
R^1f_*\mathcal O_X\simeq\mathcal O\oplus\mathcal O(-1),\qquad
\operatorname{Pic}(X)\simeq\mathbb C,
\]
and the finite-filling Hodge-line calculation now distinguishes the descended
canonical section from the direct-image frame.  In the notation of
\cref{lem:hodge-line-audit},
\[
 \operatorname{div}(\overline\sigma_j)=k_jS_j,
 \qquad
 f_{j*}\omega_{N_j/D_j}
 =\mathcal O_{D_j}\cdot s_j^{k_j}\vartheta'_j,
 \qquad k_j=m_j-1-a_j.
\]
Thus \(\overline\sigma_2\) has divisor \(2S_2\) and is not itself a
nowhere-vanishing frame, while the displayed invariant module still has the
claimed generator, a unit times \(s_j^{-a_j}e\), over the punctured cover.
This removes a local wording defect without changing the Hodge line or any
subsequent direct-image calculation.  The workbench also proves
degree-one base change at every scheme-theoretic fibre, and
\(\operatorname{ord}(\mathcal O_X(S_j)|_{S_j})=m_j\).  The smooth-fibre
period condition is exactly \(c\tau+d\mu-b=0\); hence every parameter
outside \(\mathbb R\) and a countable set gives one line bundle non-torsion
on every full fibre and on both \(S_j\).  The cusp extension is proved
intrinsically through nearby cycles and does not depend on the separately
reconstructed cellular collapse; see
\cref{thm:structure-direct-images-audit,thm:uniform-nontorsion-audit}.

\item \Verified\ The alternate Mayer--Vietoris route is now closed.  The
finite-band inclusions are exactly
$\alpha_q=(\pi_{1*},-\pi_{2*})$ in named primitive bases; the second cusp
attachment has degree-one determinant $-p$ and reduced degree-two block
$\operatorname{diag}(-p,1)$; and the dual obstruction squares commute with
the three verified Leray transgressions.  For $p=-1$ both global blocks are
 unimodular, giving an independent integral sphere certificate; see
 \cref{prop:actual-finite-band-maps,prop:cusp-mv-blocks,%
 thm:mv-leray-comparison,cor:candidate-mv-sphere}.

\item \Verified\ The Hodge diamond and Fr\"olicher sequence are now exact.
Here are the typed logarithmic objects used in the calculation:
\[
 B=\mathbb P^1,\qquad \Sigma=p_0+p_1+p_2,\qquad
 D=f^{-1}(\Sigma)_{\rm red}=W+S_1+S_2,
\]
\[
 \mathcal L=\Omega_B^1(\log\Sigma)\simeq\mathcal O_B(1),\qquad
 \mathcal A=\Omega_X^1(\log D),\qquad
 \mathcal E=\mathcal A/f^*\mathcal L,\qquad
 \mathcal G=R^1f_*\mathcal O_X.
\]
They occur in the locally free logarithmic conormal sequence
\[
 0\longrightarrow f^*\mathcal L\longrightarrow\mathcal A
 \longrightarrow\mathcal E\longrightarrow0,
\]
whose connecting maps have the exact types
\[
 \kappa_0:f_*\mathcal E\longrightarrow\mathcal L\otimes\mathcal G,
 \qquad
 \kappa_1:R^1f_*\mathcal E\longrightarrow
 \mathcal L\otimes R^2f_*\mathcal O_X\simeq\mathcal O_B.
\]
The logarithmic relative direct images are
\[
 f_*\mathcal E\simeq\mathcal O\oplus\mathcal O(-1),\qquad
 R^1f_*\mathcal E\simeq
 \mathcal O^{\oplus2}\oplus\mathcal O(-1)^{\oplus2},\qquad
 R^2f_*\mathcal E\simeq\mathcal O\oplus\mathcal O(-1),
\]
and are locally free and compatible with base change.

On a finite cover \(t=s^m\), with \((m,a)=(3,2)\) or \((4,1)\), choose
equivariant relative-cotangent frames \((e_0,e_1)\) and structure-sheaf
frames \((b_0,b_1)\) with semilinear characters
\((1,\zeta^a)\) and \((1,\zeta^{-a})\), respectively.  The actual
\(\mathbb C\{t\}\)-bases are
\[
\begin{aligned}
 f_*\mathcal E:&\quad e_0,\ s^{m-a}e_1,\\
 R^1f_*\mathcal E:&\quad
 e_0b_0,\ s^ae_0b_1,\ s^{m-a}e_1b_0,\ e_1b_1,\\
 R^2f_*\mathcal E:&\quad
 s^ae_0(b_0\wedge b_1),\ e_1(b_0\wedge b_1).
\end{aligned}
\]
Let \(\ell=dt/t\), and let
\(D_{\mathrm{KS}}=(d_{rs})_{0\le r,s\le1}\) be the marked
Kodaira--Spencer matrix after its common factor \(ds\) is removed.  Put
\[
 A=\frac{s}{m}d_{00},\qquad B_1=\frac{s^{1-a}}m d_{01},\qquad
 C=\frac{s^{1-a}}m d_{10},\qquad
 D_0=\frac{s^{m+1-2a}}m d_{11}.
\]
In source basis \((e_0,\allowbreak s^{m-a}e_1)\), target basis
\((b_0,\allowbreak s^ab_1)\ell\), ordered degree-one source basis
\((e_0b_0,\allowbreak s^ae_0b_1,\allowbreak
s^{m-a}e_1b_0,\allowbreak e_1b_1)\), and target generator
\(s^a(b_0\wedge b_1)\ell\), the coarse maps are exactly
\[
 [\kappa_0]=\begin{pmatrix}A&tB_1\\C&D_0\end{pmatrix},\qquad
 [\kappa_1]=(-C,A,-D_0,B_1).
\]
The finite residue is the typed map
\[
 \operatorname{res}_{1,j}:
 (\mathcal L\otimes\mathcal G)|_{p_j}
 \longrightarrow H^1(S_j,\mathcal O_{S_j}),\qquad
 \operatorname{res}_{1,j}\bigl((xb_0+ys^ab_1)\ell\bigr)=m_jx[b_0]
 \quad(x,y\in\mathbb C),
\]
and the derived residue square gives
\[
 \operatorname{res}_{1,j}\kappa_0=0,\qquad
 \operatorname{im}(\kappa_0\bmod t)
 =\ker\operatorname{res}_{1,j}=\mathbb C\,s^ab_1.
\]

At the cusp put
\[
 \mathcal Q_{\rm cusp}:=
 \nu_*\mathcal O_{\mathrm{dP}_6}/\mathcal O_W,
\]
and let
\(\rho:\mathcal E|_W\to\mathcal Q_{\rm cusp}\) be the relative
normalization--residue quotient map.  Retain the connecting isomorphisms
\[
 \delta_0:H^0(\mathcal Q_{\rm cusp})\xrightarrow{\ \sim\ }
 H^1(W,\mathcal O_W),\qquad
 \delta_1:H^1(\mathcal Q_{\rm cusp})\xrightarrow{\ \sim\ }
 H^2(W,\mathcal O_W).
\]
For
\[
 q_1=(-1,0,1),\qquad q_2=(0,-1,-1),\qquad
 \eta_i=d\log x_i,\qquad \alpha_i=\delta_0(q_i),
\]
let \(\theta\in H^1(\mathcal Q_{\rm cusp})\) be the unique class satisfying
\[
 \delta_1(\theta)=-\alpha_1\smile\alpha_2.
\]
The standard
\v Cech differential
\((\delta c)_{i_0\ldots i_{r+1}}
=\sum_k(-1)^kc_{i_0\ldots\widehat{i_k}\ldots i_{r+1}}\)
fixes the sign
\[
 \delta_1\bigl(\rho_*(\alpha\otimes\eta_i)\bigr)
 =-\alpha\smile\alpha_i.
\]
Consequently the two cusp maps, with their domains and codomains, are
\[
 H^0(\mathcal E|_W)\xrightarrow{\ \sim\ }
 H^0(\mathcal Q_{\rm cusp}),
\]
\[
 H^1(\mathcal E|_W)=
 \mathbb C\langle\alpha_1\eta_1,\alpha_1\eta_2,
 \alpha_2\eta_1,\alpha_2\eta_2\rangle
 \xrightarrow{\ (0,1,-1,0)\ }
 H^1(\mathcal Q_{\rm cusp})=\mathbb C\theta;
\]
the second is surjective.

Finally, the invariant section
\(\xi=\partial/\partial\zeta_2\), nowhere zero in the locally free relative
sheaf $T_{X/B}=\mathcal E^\vee$, gives the split evaluation sequence
\[
 0\longrightarrow\mathcal O(-1)\longrightarrow f_*\mathcal E
 \xrightarrow{\operatorname{ev}_\xi}\mathcal O\longrightarrow0.
\]
There is a unique lift \(h\in H^0(B,f_*\mathcal E)\) of
\(1\in H^0(B,\mathcal O)\), because
\(H^0(B,\mathcal O(-1))=0\).  On the marked cover it has the uniquely
determined form
\[
 h=(a(z),1)\binom{d\zeta_1}{d\zeta_2},\qquad a(z_2)=\frac12.
\]
Use the genuine invariant target frame
\[
 \omega_2=s\,u_2(s)(\gamma\wedge u)\ell,\qquad u_2(0)=1.
\]
The determinant and final global scalar are
\[
 \operatorname{div}(\det\kappa_0)=p_1+p_2,
 \qquad
 \kappa_1(h\otimes\gamma)|_{p_2}
 =-\frac14\dot\mu(z_2)\,\omega_2|_{p_2}\ne0
 \quad\text{in }(\mathcal L\otimes R^2f_*\mathcal O_X)|_{p_2}.
\]
This uses
\(\dot\tau(z_2)=0\) and \(K(z_2)=6\dot\mu(z_2)^2\ne0\).

Thus, with \(\mathcal K=\ker\kappa_1\),
\[
 \mathcal K\simeq\mathcal O\oplus\mathcal O(-1)^{\oplus2},
\]
and the torsion-free absolute cotangent direct images are
\[
 R^0f_*\Omega_X^1\simeq\mathcal O(-2),\quad
 R^1f_*\Omega_X^1\simeq\mathcal K,
 \quad R^2f_*\Omega_X^1\simeq\mathcal O\oplus\mathcal O(-1).
\]
Therefore
\[
 h^{1,2}=h^{2,1}=1,\qquad h^{1,1}=h^{2,2}=2.
\]
In full, with row index \(q\) and column index \(p\), the Hodge table is
\[
\begin{array}{c|cccc}
q\backslash p&0&1&2&3\\ \hline
3&0&0&0&1\\
2&0&1&2&1\\
1&1&2&1&0\\
0&1&0&0&0.
\end{array}
\]
The two middle \(d_1\)-rows are exact
\(0\to\mathbb C\to\mathbb C^2\to\mathbb C\to0\), and
\[
 \dim_{\mathbb C}E_1=10,\qquad
 \dim_{\mathbb C}E_2=2,\qquad E_2=E_\infty.
\]
See
\cref{prop:log-relative-direct-images,prop:log-ks-divisor,%
prop:absolute-cotangent-direct-images,thm:h12-scalar,cor:exact-hodge-diamond}.

\item \Refuted\ Remark~1.3's assertion that the relative differential
sheaves have torsion along all three singular fibres is false already for
$\Omega^1_{X/B}$.  The $t=xy$ and $t=xyz$ modules inject into the free
logarithmic relative module, so $\Omega^1_{X/B}$ is torsion-free but
non-locally-free near the cusp.  At a point of the smooth finite cover
\(q_j:J'_j\to N_j\), with
\(R_j=\mathcal O_{J'_j,x}\simeq\mathbb C\{s,y_1,y_2\}\), and
\(t=s^{m_j}\), the exact module is
\[
 (q_j^*\Omega^1_{N_j/B})_x
 \simeq\Omega^1_{J'_j/B,x}
 \simeq R_j\,dy_1\oplus R_j\,dy_2
 \oplus\bigl(R_j/(s^{m_j-1})\bigr)ds,
\]
and hence
\[
 (q_j^*\operatorname{tors}\Omega^1_{N_j/B})_x
 =\bigl(R_j/(s^{m_j-1})\bigr)ds.
\]
Faithful flatness descends this module and its support.  Thus the torsion
support is exactly the two multiple fibres and is absent at the cusp; see
the local calculation (11.17f).

\item \Verified\ The Bott--Chern and Aeppli dimensions are now exact as
well.  McHugh's threefold formula gives
\[
h^{1,2}_{BC}=h^{1,2}+\dim E_2^{3,1}-k^{2,0}=1+0-0=1,
\]
because the de Rham image of the pure doubly closed forms surjects onto
$K^{2,0}$, so $k^{2,0}\le b_2=0$.  The complete Bott--Chern rows
$q=3,2,1,0$ are
\[
(0,0,1,1),\quad(0,1,2,1),\quad(0,2,1,0),\quad(1,0,0,0),
\]
and Aeppli duality gives
\[
(0,0,0,1),\quad(0,1,2,0),\quad(1,2,1,0),\quad(1,1,0,0).
\]
See \cref{thm:exact-bott-chern-aeppli}; no dimension remains open.

\item \Verified\ Global representatives and the complete Bott--Chern
product structure are now known.  The two Dolbeault inputs must not be
conflated.  The distinguished class
\[
 \mathfrak g\in H^{0,1}_{\bar\partial}(X)=H^1(X,\mathcal O_X)
\]
is the unique coherent-Leray class mapping to the invariant section
\(\gamma^{0,1}\in H^0(B,R^1f_*\mathcal O_X)\).  In contrast, the marked
integral fibre class is
\[
 q=\eta=u\wedge w+6\gamma\wedge\delta\in H^2(F_b;\mathbb Z),
\]
whose holomorphic extension is a section
\(\boldsymbol\eta\in H^0(B,R^1f_*\Omega_X^1)\).  Choose a coherent-Leray
lift
\[
 \mathfrak q\in H^{1,1}_{\bar\partial}(X)
 \quad\text{with edge value}\quad \boldsymbol\eta.
\]
This lift is not distinguished: its exact ambiguity is
\[
 \mathfrak q\longmapsto\mathfrak q+c\,d_1\mathfrak g,
 \qquad c\in\mathbb C,
\]
because the kernel of the edge map is
\(H^1(B,f_*\Omega_X^1)=\mathbb C\,d_1\mathfrak g\).

Choose global \(\bar\partial\)-closed smooth representatives
\[
 G\in\mathcal A^{0,1}(X),\quad [G]_{\bar\partial}=\mathfrak g,
 \qquad
 Q\in\mathcal A^{1,1}(X),\quad [Q]_{\bar\partial}=\mathfrak q.
\]
For a marked smooth fibre put
\[
 y=\operatorname{Im}\tau,\qquad m_\mu=\operatorname{Im}\mu,\qquad
 r=\operatorname{Im}\beta,\qquad
 \Delta_F=r-\frac{6m_\mu^2}{y}<0,
\]
and define the translation-invariant forms
\[
 \Gamma=\frac{i}{2\Delta_F}
 \left(d\bar\zeta_2-\frac{m_\mu}{y}d\bar\zeta_1\right),
 \qquad
 U=\frac{i}{2y\Delta_F}
 \left(r\,d\bar\zeta_1-6m_\mu\,d\bar\zeta_2\right),
\]
\[
 Q_F=-d\zeta_1\wedge U-6d\zeta_2\wedge\Gamma.
\]
The restriction statements are equalities of typed Dolbeault classes,
not claims about a preferred literal coefficient gauge:
\[
 i_b^*\mathfrak g=[\Gamma]_{\bar\partial}
   \in H^{0,1}_{\bar\partial}(F_b),
 \qquad
 i_b^*\mathfrak q=[Q_F]_{\bar\partial}
   \in H^{1,1}_{\bar\partial}(F_b).
\]
After fixing \(\mathfrak g\) and one \(\mathfrak q\), changing smooth
representatives by
\(G\mapsto G+\bar\partial\phi\) and
\(Q\mapsto Q+\bar\partial\Psi\) changes the derived representatives only
by \(\partial\bar\partial\)-exact terms.  Thus no preferred literal global
coefficient atlas through the singular fillings is asserted or needed.
Put
\[
 \mathsf A:=[\partial G]_{BC},\qquad
 \mathsf B:=[\bar\partial\bar G]_{BC}
 \quad\text{in }H^{1,1}_{BC}(X).
\]

Let
\[
 \exp:H^1(X,\mathcal O_X)\longrightarrow\operatorname{Pic}(X)
\]
be induced by \(z\mapsto e^{2\pi iz}\), and orient
\(\omega\in H^2(B;\mathbb Z)\) by
\(\omega=c_1(\mathcal O_B(1))\).  The integral Picard normalization and
Bott--Chern first Chern class are
\[
 \exp(12\mathfrak g)=f^*\mathcal O_B(1),\qquad
 c_{1,BC}(\exp(z\mathfrak g))=z\mathsf A+\bar z\,\mathsf B,
\]
\[
 c_{1,BC}(f^*\mathcal O_B(1))=12(\mathsf A+\mathsf B).
\]
The absolute factor is \(12\); choosing the opposite primitive orientation
replaces \(\mathfrak g\) by \(-\mathfrak g\) in the pullback identity.
The exact fibre pairings, with the orientation stated, are
\[
 \int_{(F_b,\operatorname{or}_{\rm mark})}Q_F^2=12,
 \qquad
 \int_{(F_b,\operatorname{or}_{\mathbb C})}Q_F^2=-12,
\]
\[
 \int_{(F_b,\operatorname{or}_{\mathbb C})}
 Q_F\wedge\Gamma\wedge d\zeta_2=-1.
\]
They prove the needed Dolbeault nonvanishing, while the constant fibre
coefficients give the typed vanishing
\[
 i_b^*d_1[G\wedge Q]_{\bar\partial}
 =i_b^*[\partial(G\wedge Q)]_{\bar\partial}=0
 \quad\text{in }H^{2,2}_{\bar\partial}(F_b).
\]
Choose \(\upsilon\in\mathcal A^{3,3}(X)\) with \(\int_X\upsilon=1\).
Together with the unit
\(\mathbf1=[1]_{BC}\in H^{0,0}_{BC}(X)\), define
\[
\begin{array}{c|c}
\text{class}&\text{bidegree}\\ \hline
\mathsf A=[\partial G]_{BC},\quad
 \mathsf B=[\bar\partial\bar G]_{BC}&(1,1)\\
\mathsf C=[\partial Q]_{BC},\quad
 \overline{\mathsf C}=[\bar\partial\bar Q]_{BC}&(2,1),(1,2)\\
\mathsf D=[\partial(G\wedge Q)]_{BC},\quad
 \overline{\mathsf D}=[\bar\partial(\bar G\wedge\bar Q)]_{BC}
 &(2,2),(2,2)\\
\mathsf E=[\partial(Q^2)]_{BC},\quad
 \overline{\mathsf E}=[\bar\partial(\bar Q^2)]_{BC}
 &(3,2),(2,3)\\
\mathsf T=[\upsilon]_{BC},\quad\int_X\mathsf T=1&(3,3).
\end{array}
\]
These are bases of every nonzero positive-degree Bott--Chern group.  For
their span \(N_{\mathrm{BC}}\),
\[
 H^{*,*}_{BC}(X)=\mathbb C\mathbf1\oplus N_{\mathrm{BC}},\qquad
 N_{\mathrm{BC}}^2=0.
\]
See
\cref{prop:bcpr-generators-are-bases,thm:bcpr-square-zero-algebra}.  The
lift and smooth-gauge ambiguities above change neither the verified
dimensions nor the square-zero product algebra.

\item \Verified\ The positive radial cusp now has an explicitly declared
flat quotient metric and an exact scalar eigenvalue calculation.  Write
\[
 \mu_T=\alpha_T+i m_T,\qquad h_T=\eta_T+i H_T,\qquad
 b_T=\xi_T+i B_T,\qquad c_T=\xi_T-\eta_T,
\]
\[
 L_T=T+H_T,\qquad q_T=L_T-B_T,\qquad
 D_T=L_Tq_T+6m_T^2.
\]
Set also
\[
 m=\operatorname{Im}\mu(0),\qquad
 H=\operatorname{Im}h(0),\qquad
 B=\operatorname{Im}b(0).
\]
In the ordered lattice basis
\((\widehat\gamma,\widehat u,\widehat w,\widehat\delta)\)
and real coordinate order \((x_1,y_1,x_2,y_2)\), the real period
matrix, its oriented determinant, and its positive covolume are
\[
 P_T=
 \begin{pmatrix}
  6\alpha_T&\eta_T&1&0\\
  6m_T&L_T&0&0\\
  c_T&\alpha_T&0&1\\
  -q_T&m_T&0&0
 \end{pmatrix},\qquad
 \det P_T=-D_T,\qquad
 \operatorname{covol}(P_T\mathbb Z^4)=D_T>0,\qquad
 D_{\mathrm{can},T}=-\frac{D_T}{L_T}<0.
\]
The full inverse-dual certificate is
\[
 P_T^{-\mathsf T}=
 \begin{pmatrix}
 0&0&1&0\\[2pt]
 \dfrac{m_T}{D_T}&\dfrac{q_T}{D_T}&
 \dfrac{-6m_T\alpha_T-q_T\eta_T}{D_T}&
 \dfrac{-m_Tc_T-q_T\alpha_T}{D_T}\\[7pt]
 0&0&0&1\\[2pt]
 -\dfrac{L_T}{D_T}&\dfrac{6m_T}{D_T}&
 \dfrac{6(L_T\alpha_T-m_T\eta_T)}{D_T}&
 \dfrac{L_Tc_T-6m_T\alpha_T}{D_T}
 \end{pmatrix}.
\]
Thus \((P_T\mathbb Z^4)^*=P_T^{-\mathsf T}\mathbb Z^4\), and for
\(k=(r,s,n_1,n_2)^{\mathsf t}\), with
\[
 R_T(k)=r-6\alpha_Tn_1-c_Tn_2,\qquad
 S_T(k)=s-\eta_Tn_1-\alpha_Tn_2,
\]
its member indexed by \(k\) is exactly
\[
 \left(n_1,
 \frac{m_TR_T(k)+q_TS_T(k)}{D_T},n_2,
 \frac{-L_TR_T(k)+6m_TS_T(k)}{D_T}\right).
\]
For the declared Euclidean quotient metric \(g_T\), the first positive
scalar eigenvalue is, for all large \(T\),
\[
 \lambda_1(F_T,g_T)
 =4\pi^2\min\left\{
 \frac{L_T^2+m_T^2}{D_T^2},
 \frac{q_T^2+36m_T^2}{D_T^2}\right\}
 =\frac{4\pi^2}{T^2}\left(
 1+\frac{-2H+2\min\{B,0\}}{T}+O(T^{-2})\right).
\]
The metric normalization is recorded separately and gives
\[
 \begin{aligned}
 \widehat g_T&=D_T^{-1/2}g_T,&
 \operatorname{Vol}(F_T,\widehat g_T)&=1,\\
 \lambda_1(F_T,\widehat g_T)
 &=\frac{4\pi^2}{T}\left(
  1+\frac{-H-B/2+2\min\{B,0\}}{T}+O(T^{-2})\right),\\
 \lambda_1(F_T,\widehat g_T)&\longrightarrow0.
 \end{aligned}
\]
Normalized Haar measure on the fixed angular torus is only the unitary
measure convention used below; it is not the metric rescaling
\(g_T\mapsto\widehat g_T\).  See
 \cref{cusp-spec-exact-dual-lattice,cusp-spec-shortest-modes,%
 cusp-spec-unit-volume}.

\item \Verified\ The degree-
\(N^2\) regular-fibre cover extends across the completed local toric cusp,
with the base-change and total degrees kept distinct.  Retain
\[
 B_0=\begin{pmatrix}0&1\\-1&0\end{pmatrix},\qquad
 Z(s)=sB_0+C(t),\qquad t=e^{2\pi i s},
\]
let \(Y=Y_\Sigma\) be the toric threefold of the height-one
translation-invariant \(A_2\) fan, and let the target action be
\[
 \Psi_{\bar\lambda}(x,t)
 =\bigl(e^{2\pi iC(t)\bar\lambda}
        t^{B_0\bar\lambda}x,t\bigr),\qquad
 f_0:N_0:=Y/\bar\Lambda\longrightarrow\Delta_t.
\]
For every integer \(N\ge1\), after shrinking the disc, put
\[
 \rho_N(v)=v^N,\qquad
 C_N(v)=\frac1N C(v^N),\qquad
 \Psi^{(N)}_{\bar\lambda}(\xi,v)
 =\bigl(e^{2\pi iC_N(v)\bar\lambda}
        v^{B_0\bar\lambda}\xi,v\bigr).
\]
This is a free properly discontinuous action on a copy \(Y^{(N)}\) of
\(Y\), and
\[
 f_N:N_0^{\langle N\rangle}:=Y^{(N)}/\bar\Lambda
 \longrightarrow\Delta_v
\]
is a proper holomorphic map from a smooth threefold, with reduced
irreducible nonnormal central fibre \(W_N\) normalized by
\(\mathrm{dP}_6\).  Multiplication by
\(N\) on the three-dimensional cocharacter lattice gives
\[
 P_N:Y^{(N)}\longrightarrow Y,\qquad
 (\xi_1,\xi_2,v)\longmapsto(\xi_1^N,\xi_2^N,v^N),
\]
whose maximal-chart ring map is
\[
 \mathbb C\{z_0,z_1,z_2\}\longrightarrow
 \mathbb C\{w_0,w_1,w_2\},\qquad z_i\longmapsto w_i^N.
\]
It is finite flat of degree \(N^3\), is equivariant for the two displayed
actions, and descends to a finite flat degree-
\(N^3\) map
\[
 \kappa_N:N_0^{\langle N\rangle}\longrightarrow N_0,\qquad
 f_0\circ\kappa_N=\rho_N\circ f_N,\qquad
 \kappa_N^*W=NW_N,\qquad
 \operatorname{Ram}(\kappa_N)=(N-1)W_N.
\]
For the normal base change
\(N_0[N]=N_0\times_{\Delta_t}\Delta_v\), the unique factorization
\[
 N_0^{\langle N\rangle}\xrightarrow{Q_N}N_0[N]\longrightarrow N_0
\]
has generic degree \(N^2\).  If \(N>1\), \(Q_N\) is quasi-etale and
nonflat: its non-etale and nonflat loci are precisely the inverse images of
the double and triple strata, where the closed-fibre lengths are respectively
\(2N-1>N\) and \(3N^2-3N+1>N^2\).  On reduced central fibres the generic
degree is \(N^2\), and the normalization square is
\[
\begin{CD}
 \mathrm{dP}_6 @>{[N]}>> \mathrm{dP}_6\\
 @V{\nu_N}VV @VV{\nu}V\\
 W_N @>>> W,
\end{CD}
\]
where \([N]\) has degree \(N^2\) and ramification index \(N\) on each of
the six toric boundary curves.  Over \(v\ne0\), \(Q_N\) is the cover for
\[
 \mathbb Z\langle\widehat\gamma,\widehat u,
 N\widehat w,N\widehat\delta\rangle\subset\Lambda,
 \qquad \operatorname{Deck}(Q_N)\cong(\mathbb Z/N\mathbb Z)^2,
\]
and, with \(S_N=\operatorname{diag}(1,1,N,N)\),
\[
 S_N^{-1}M_0^NS_N=M_0.
\]
This is a local completed-cusp morphism, not a degree-
\(N^2\) self-map over the unchanged disc, not a normalization, and not a
cover over the entire base; see
\cref{cusp-cover-extension-theorem,cusp-cover-naive-action-failure,%
cusp-cover-global-scope}.
Its use in four-dimensional finite-cover unspooling is an ancillary
construction: it changes both the fibre lattice and the cusp base.  It does
not answer the native \(S^{12}\)/\(S^{24}\) sphere-route problem, whose
solution begins from the unaltered quotient fibres \(Q_T\) and \(C_T\).

\item \Verified\ The native cusp quotient and the invariant-circle quotient
give distinct exact conical smash-sphere carriers.  With
\[
 K_\Lambda=\mathbb Z\langle\widehat w,\widehat\delta\rangle
 =\operatorname{im}(M_0-I),\qquad
 \overline\Lambda=\Lambda/K_\Lambda,\qquad
 \overline T^2=(\overline\Lambda\otimes\mathbb R)/\overline\Lambda,
\]
there is a fibrewise real-analytic group epimorphism
\[
 \overline\pi:N_0\setminus W\longrightarrow D_0^*\times\overline T^2
\]
with connected kernel and trivial target local system.  On
\(F_T=f^{-1}(e^{-2\pi T})\), it is \(\pi_T\) after the isometry
\[
 \iota_T:\overline T^2\xrightarrow{\sim}Q_T,\qquad
 [(a,b)]\longmapsto[Y_T(a,b)^{\mathsf T}],\qquad
 Y_T=\begin{pmatrix}6m_T&L_T\\-q_T&m_T\end{pmatrix},
\]
 and, for a fixed constant \(C_{\mathrm{sys}}>0\),
 \(\operatorname{sys}(Q_T)\ge T-C_{\mathrm{sys}}\).  For
\[
 \mathcal A=
 \frac{[0,\epsilon]\times S^1\times\overline T^2}
 {\{0\}\times S^1\times\overline T^2\sim v_0,
  \ \{\epsilon\}\times S^1\times\overline T^2\sim v_1}
 \cong\operatorname{Susp}(T^3),
\]
the map
\[
 q_2:(X,x_0)\longrightarrow(\mathcal A,v_0),\qquad
 q_2(x)=[r,e^{i\phi},\overline\pi_f(x)]
\]
for \(t(f(x))=re^{i\phi}\), with
\(\overline\pi_f=\operatorname{pr}_{\overline T^2}\circ\overline\pi\),
\(q_2(W)=v_0\) and
\(q_2(X\setminus\operatorname{int}N_0)=v_1\), is a continuous surjective
based quotient.  Hence
\[
 q_2^{\wedge2}:X\wedge X\cong S^{12}\longrightarrow
 \mathcal A\wedge\mathcal A
\]
is global, and its synchronized radial restriction is exactly
\[
 \pi_T\times\pi_T:F_T\times F_T\longrightarrow Q_T\times Q_T,
 \qquad
 (Q_T\times Q_T,0)
 \xrightarrow[T\to\infty]{\mathrm{pointed}}(\mathbb R^4,0)
\]
in pointed local-isometric, pointed smooth, and pointed
Gromov--Hausdorff topology.

The cusp-supported one-circle route uses
\[
 \mathcal B=
 \frac{[0,\epsilon]\times S^1\times S^1}
 {\{0\}\times S^1\times S^1\sim b_0,\ 
  \{\epsilon\}\times S^1\times S^1\sim b_1}
 \cong\operatorname{Susp}(T^2)
\]
and the continuous surjective based quotient
\[
 q_{1,\mathrm{cusp}}:(X,x_0)\longrightarrow(\mathcal B,b_0),\qquad
 q_{1,\mathrm{cusp}}(x)=[r,e^{i\phi},\chi_\gamma(x)]
\]
on \(N_0\setminus W\), with \(q_{1,\mathrm{cusp}}(W)=b_0\) and the
exterior collapsed to \(b_1\).  Its fourth smash power is
\[
 q_{1,\mathrm{cusp}}^{\wedge4}:X^{\wedge4}\cong S^{24}
 \longrightarrow\mathcal B^{\wedge4},
\]
and its synchronized radial restriction is
\(\chi_{\gamma,T}^{\times4}:F_T^4\to C_T^4\).

For the full-base circle route, let
\[
 P=B\setminus\bigcup_{j=0}^2\operatorname{int}D_j,\qquad
 \mathcal M=(P\times S^1)\underset{h_1^{S^1}}\cup C_1
                         \underset{h_2^{S^1}}\cup C_2,\qquad
 \partial\mathcal M=\partial D_0\times S^1,
\]
and
\(\mathcal C=\mathcal M\cup_{\partial\mathcal M}
\operatorname{Cone}(\partial\mathcal M)\), with cone vertex \(c_0\).
The continuous surjective based quotient
\[
 q_1:(X,x_0)\longrightarrow(\mathcal C,c_0)
\]
 equals \((f,\theta_\gamma)\) over \(P\), equivalently
 \(\chi_\gamma|_{f^{-1}(P)}\), and equals the finite-filling maps
\(N_j\to C_j\) for \(j=1,2\), sends
\(x\in N_0\setminus W\) to
 \([r,(e^{i\phi},\theta_\gamma(x))]\), and sends \(W\) to \(c_0\).  Its
fourth smash power is
\[
 q_1^{\wedge4}:X^{\wedge4}\cong S^{24}\longrightarrow
 \mathcal C^{\wedge4},
\]
whose synchronized radial restriction is exactly
\[
\begin{aligned}
 \chi_{\gamma,T}^{\times4}&:F_T^4\longrightarrow C_T^4,\\
 C_T&=\mathbb R/\ell_T\mathbb Z,\\
 \ell_T&=\frac{D_T}{\sqrt{L_T^2+m_T^2}}
       =T+H-B+O(T^{-1}),
\end{aligned}
\]
so
\[
 (C_T^4,0)
 \xrightarrow[T\to\infty]{\mathrm{pointed}}(\mathbb R^4,0)
\]
in the same three pointed topologies.

The exponents \(2\) and \(4\), and hence the source-sphere dimensions
\(12\) and \(24\), are optimal for these two native quotient problems with
the unscaled quotient of the marked Euclidean fibre metric.  Indeed, if a
primitive kernel \(N_{\mathrm{ker}}\subset\Lambda^{\oplus r}\) defines a
connected torus
quotient of \(r\) cusp fibres whose quotient systole tends to infinity, then
every uniformly bounded period must lie in \(N_{\mathrm{ker}}\).  Thus
\[
 K_\Lambda^{\oplus r}\subset N_{\mathrm{ker}},\qquad
 \operatorname{rank}(\Lambda^{\oplus r}/N_{\mathrm{ker}})\le 2r.
\]
For a quotient over the full regular base, or over any finite cover of that
base, monodromy forces the additional sublattice
\(\mathbb Z\langle\widehat u\rangle^{\oplus r}\) into
\(N_{\mathrm{ker}}\): one uses
\(A_1\widehat w=\widehat u-\widehat w\), and on a finite cover the exact
formula
\[
 (A_2A_1)^m\widehat w
 =\widehat w+m(\widehat u-\widehat\delta)
\]
together with primitivity of \(N_{\mathrm{ker}}\).  Consequently
\[
 \operatorname{rank}(\Lambda^{\oplus r}/N_{\mathrm{ker}})\le r.
\]
A rank-four expanding quotient therefore needs at least two copies
cusp-locally and at least four copies globally; the displayed
\(Q_T^2\) and \(C_T^4\) constructions attain the two bounds.

After collapsing the coordinate one-skeleton
\[
 K^{(1)}=
 ((\mathbb R/\mathbb Z)\times\{0\})
 \cup(\{0\}\times(\mathbb R/\mathbb Z))
 \subset\overline T^2,
\]
the two native routes also give continuous surjective based quotient maps
\[
 \widehat q_3:X\twoheadrightarrow S^3,\qquad
 \widehat q_\gamma:X\twoheadrightarrow S^2,
\]
and hence, using the fixed based choice
\(X\cong_{\mathrm{diff}}S^6\),
\[
 \widehat q_3^{\wedge2}:S^{12}\twoheadrightarrow S^6,\qquad
 \widehat q_\gamma^{\wedge4}:S^{24}\twoheadrightarrow S^8.
\]
Both full sphere carriers are based null-homotopic.  First let
\(f=\widehat q_\gamma:S^6\to S^2\) after the chosen based identification.
There is an exact factorization
\[
 S^{12}=S^6\wedge S^6
 \xrightarrow{\ f\wedge\operatorname{id}_{S^6}=\Sigma^6f\ }
 S^2\wedge S^6=S^8
 \xrightarrow{\ \operatorname{id}_{S^2}\wedge f\ }
 S^2\wedge S^2=S^4.
\]
Freudenthal stability gives
\(\pi_{12}(S^8)\cong\pi_4^{\mathrm S}=0\), so \(f\wedge f\) is based
null-homotopic.  Since
\(f^{\wedge r}=(f\wedge f)\wedge f^{\wedge(r-2)}\), it follows that
\[
 [f^{\wedge r}]=0\in\pi_{6r}(S^{2r})
 \qquad\text{for every }r\ge2.
\]
In particular, the displayed surjective quotient
\(\widehat q_\gamma^{\wedge4}:S^{24}\twoheadrightarrow S^8\) is itself
based null-homotopic.

For the other carrier, let \(g:S^6\to S^3\) be any based map.  Its exact
smash factorization is
\[
 g\wedge g=(\Sigma^3g)\circ(\Sigma^6g):
 S^{12}\longrightarrow S^9\longrightarrow S^6.
\]
Write \(\operatorname{st}\) for stabilization.  The exact low-stem data are
\[
 \pi_6(S^3)\cong\mathbb Z/12,
 \quad
 \pi_3^{\mathrm S}=\mathbb Z/24\langle\nu\rangle,
 \quad
 \pi_6^{\mathrm S}=\mathbb Z/2\langle\nu^2\rangle,
\]
\[
 \pi_{12}(S^6)=\mathbb Z/2\langle\nu_6\circ\nu_9\rangle,
 \qquad
 \operatorname{st}:\pi_{12}(S^6)
 \xrightarrow{\ \cong\ }\pi_6^{\mathrm S}.
\]
The full groups are tabulated in \cite[p.~339]{Hatcher2002}; the signed smash
formula and the unstable stabilization calculation are
\cite[Propositions~3.1 and~5.11]{Toda1962}.
Since \(12[g]=0\), one has
\(\operatorname{st}([g])=2a\nu\) for some \(a\in\mathbb Z/12\).  By
multiplicativity,
\[
 \operatorname{st}([g\wedge g])
 =\operatorname{st}([g])^2
 =4a^2\nu^2=0\in\pi_6^{\mathrm S}.
\]
The stabilization isomorphism forces
\([g\wedge g]=0\in\pi_{12}(S^6)\).
Thus \(g^{\wedge r}\) is based null-homotopic for every \(r\ge2\).  Taking
\(g=\widehat q_3\) proves
\[
 [\widehat q_3^{\wedge2}]=0\in\pi_{12}(S^6),
 \qquad
 [\widehat q_\gamma^{\wedge4}]=0\in\pi_{24}(S^8).
\]
Both displayed maps remain surjective quotient maps.  The one-copy classes
\([\widehat q_3]\in\pi_6(S^3)\) and
\([\widehat q_\gamma]\in\pi_6(S^2)\) have not been computed.
Their synchronized radial restrictions contain based topologically embedded
four-spheres, equipped separately with the intrinsic quotient length metrics
\[
 \mathbb S_{12,T}^4
   =(Q_T/K_T^{(1)})^{\wedge2}\hookrightarrow S^6,\qquad
 \mathbb S_{24,T}^4=C_T^{\wedge4}\hookrightarrow S^8,
\]
obtained from, respectively,
\[
 F_T^2\xrightarrow{\ \pi_T^{\times2}\ }Q_T^2
 \xrightarrow{\ c_T^{\times2}\ }(S_T^2)^2
 \longrightarrow\mathbb S_{12,T}^4,\qquad
 F_T^4\xrightarrow{\ \chi_{\gamma,T}^{\times4}\ }C_T^4
 \longrightarrow\mathbb S_{24,T}^4,
\]
where the last arrows collapse the fat wedges.  The full maps from
\(F_T^2\) and \(F_T^4\) retain the nontrivial kernels of their first arrows;
only the quotient-stage maps are injective off their explicitly collapsed
skeleton/fat-wedge loci.  Equip the displayed products
with their Euclidean \(\ell^2\)-product metrics and give each four-sphere the
induced quotient length metric.  If
\(z_T\) is the image of \([(1/2,1/2)]\) and
\[
 \zeta_{12,T}=[z_T,z_T],\qquad
 \zeta_{24,T}=
 [\ell_T/2,\ell_T/2,\ell_T/2,\ell_T/2],
\]
then every radius-\(R\) ball about these points is exactly a Euclidean
four-ball once, respectively,
\[
 T-C_{\mathrm{sys}}>4R,\qquad \ell_T>4R.
\]
Therefore
\[
 (\mathbb S_{12,T}^4,\zeta_{12,T})\longrightarrow(\mathbb R^4,0),\qquad
 (\mathbb S_{24,T}^4,\zeta_{24,T})\longrightarrow(\mathbb R^4,0)
\]
in pointed local-isometric and pointed Gromov--Hausdorff topology, with smooth
convergence on the eventual exact Euclidean neighborhoods of the displayed
base points.  Both based inclusion maps
\[
 i_{12,T}:\mathbb S_{12,T}^4\hookrightarrow S^6,
 \qquad
 i_{24,T}:\mathbb S_{24,T}^4\hookrightarrow S^8
\]
are null-homotopic: sliding their fixed interior suspension coordinate to a
suspension endpoint contracts them through based maps.  Thus
\((i_{12,T})_*[S^4]=0\in H_4(S^6;\mathbb Z)\) and
\((i_{24,T})_*[S^4]=0\in H_4(S^8;\mathbb Z)\).  These slice-inclusion
null-homotopies are distinct from the global equalities
\([\widehat q_3^{\wedge2}]=0\in\pi_{12}(S^6)\) and
\([\widehat q_\gamma^{\wedge4}]=0\in\pi_{24}(S^8)\).

Deleting the smash basepoint nevertheless gives literal nonclosed
topologically embedded copies of \(\mathbb R^4\):
\[
 (Q_T\setminus K_T^{(1)})^2
 \cong\mathbb S_{12,T}^4\setminus\{*\}
 \cong\mathbb R^4\hookrightarrow S^6,
\qquad
 (C_T\setminus\{0\})^4
 \cong\mathbb S_{24,T}^4\setminus\{*\}
 \cong\mathbb R^4\hookrightarrow S^8.
\]
They are open only in the embedded four-spheres and are not globally
Euclidean-isometric under a fixed compatible compact-target metric.

There are also Hausdorff, second-countable topological five-manifolds with
boundary \(\mathfrak E_{12}\) and \(\mathfrak E_{24}\) whose positive fibres
are \(Q_{1/s}^{\times2}\) and \(C_{1/s}^{\times4}\), respectively, and whose
zero fibres are literally \(\mathbb R^4\).  Continuous surjections
\(\beta_\nu:\mathfrak E_\nu\to\mathcal D_\nu\) identify every positive fibre
with its radial carrier slice and collapse the zero fibre to the cone vertex.
They are homeomorphisms off the zero fibre but are nonproper and nonquotient;
thus they are comparison/collapse maps, not topological blow-downs or
resolutions, and \(\mathcal D_\nu\) is not the quotient
\(\mathfrak E_\nu/\mathbb R^4\) via \(\beta_\nu\).

The metric quantifiers are essential.  For every fixed metric compatible
with the compact conical target topology, the target slices
\(j_T(Q_T^2)\), \(k_T^{\mathcal B}(C_T^4)\), and
\(k_T^{\mathcal C}(C_T^4)\) converge in Hausdorff, hence
Gromov--Hausdorff, distance to the smash basepoint.  A fixed cusp-adapted
metric can realize the intrinsic limits only after a moving-point blow-up:
with \(r_T=e^{-2\pi T}\), the factors \(T/r_T\) and
\(\ell_T/r_T\) give the exact induced slice metrics
\[
 \left(\frac{T}{r_T}\right)^2j_T^*g_{Z_{12}}
   =g_{Q_T}\oplus g_{Q_T},\qquad
 \left(\frac{\ell_T}{r_T}\right)^2
   (k_T^{\mathcal D})^*g_{Z_{24}^{\mathcal D}}
   =g_{C_T}^{\oplus4},\quad
 \mathcal D\in\{\mathcal B,\mathcal C\}.
\]
The slice bubbles are \(\mathbb R^4\), but the corresponding full ambient
bubbles are \(\mathbb R^8\) for the \(S^{12}\) route and
\(\mathbb R^{12}\) for either \(S^{24}\) route.  More generally, if
\(\lambda_T>0\) and \(\beta_T=\lambda_Tr_T\), the slice limit is a point
when \(\beta_T\to0\), a square flat four-torus of side \(L\) when
\(\beta_T\to L\in(0,\infty)\), and \(\mathbb R^4\) when
\(\beta_T\to\infty\); in the last case the cone vertex is at scaled
distance \(\beta_T\to\infty\).  Hence the \(\mathbb R^4\) limits are
neither vertex tangent cones nor ambient limits of fixed smooth metrics on
the source spheres, whose ordinary tangent dimensions are \(12\) and \(24\).

The original maps \(q_2^{\wedge2}\), \(q_{1,\mathrm{cusp}}^{\wedge4}\),
and \(q_1^{\wedge4}\) have the displayed conically stratified codomains.
The repeated-smash null theorems for the sphere-valued composites
\(\widehat q_3\) and \(\widehat q_\gamma\) do not supply null-homotopies
before the final one-skeleton/base-angle quotients.  The unstable classes of
\[
 q_2^{\wedge2}:S^{12}\to\mathcal A^{\wedge2},
 \qquad
 q_1^{\wedge4}:S^{24}\to\mathcal C^{\wedge4}
\]
remain \Open: their zero ordinary Hurewicz images and absence of homotopy
sections do not decide nullity.  The class of
\(q_{1,\mathrm{cusp}}^{\wedge4}:S^{24}\to\mathcal B^{\wedge4}\) has
likewise not been computed.
The sphere codomains occur only after the additional one-skeleton and
base-angle collapses defining \(\widehat q_3\) and
\(\widehat q_\gamma\).  All identifications of the sources with
\(S^{12}\) and \(S^{24}\) use the chosen based diffeomorphism
\(X\cong_{\mathrm{diff}}S^6\).  See
\cref{cusp-cone-native-descent,cusp-cone-S12-route,%
cusp-cone-S24-local-route,cusp-cone-S24-route,cusp-cone-optimal-exponents,%
cusp-cone-angle-forgotten-sphere-carriers,cusp-cone-S12-sphere-map-null,%
cusp-cone-S24-sphere-map-null,%
cusp-cone-sphere-map-homotopy-status,cusp-cone-moving-S4-null-homotopy,%
cusp-cone-literal-R4-open-cells,%
cusp-cone-exact-decompactification-space,%
sphere-metric-uniform-collapse,sphere-metric-fixed-realization,%
sphere-metric-scaling-trichotomy,sphere-metric-tangent-scope,%
cusp-cone-metric-scope,cusp-cone-target-homotopy-types,%
cusp-cone-obstructions}.

\item \Refuted\ The assertion that the literal fixed-target map
\[
 \chi_\gamma:J\longrightarrow B^\circ\times(\mathbb R/\mathbb Z)
\]
extends continuously across \(W\), or that
\[
 \chi_\gamma^{[4]}:J^{[4]}\longrightarrow
 B^\circ\times(\mathbb R/\mathbb Z)^4
\]
extends with the stated regular-fibre restriction to the compact fourfold
fibre product or to \(X^4\), is false: two explicit paths have the same
toric limit in \(W\) and constant character values \(3/4\) and \(1/2\).
The cone map above changes the codomain by collapsing
\(\partial D_0\times S^1\) to \(c_0\), so it does not contradict this
nonextension.  Moreover
\[
 X^2\simeq S^6\times S^6\not\simeq S^{12},\qquad
 X^4\simeq(S^6)^4\not\simeq S^{24};
\]
only the smash products satisfy
\(X^{\wedge2}\cong S^{12}\) and \(X^{\wedge4}\cong S^{24}\).
Neither smash sphere admits an almost-complex structure, and neither route
constructs a compact complex sphere; see
\cref{cusp-spec-globalization-obstruction,cusp-cone-obstructions}.

\item \Verified\ The exceptional first-coordinate elliptic quotient is now
classified on the actual period image, rather than on an arbitrary ambient
path.  For
\[
 \mathcal A_z=\mathbb C^2/
 \mathbb Z\langle(6\mu,\beta),(\tau,\mu),(1,0),(0,1)\rangle,
\]
the first coordinate induces a surjective holomorphic homomorphism
\[
 \mathcal A_z\longrightarrow\mathbb C/(\mathbb Z\tau+\mathbb Z)
\]
if and only if \(6\mu=a\tau+b\) for a unique
\((a,b)\in\mathbb Z^2\).  Its connected kernel has period
\(\tau_f=a\mu-\beta\) and
\(\operatorname{Im}\tau_f=-D_{\mathrm{can}}>0\).  The exact marked-label
action is
\[
 g_0:(a,b)\mapsto(a,a+b),\qquad
 g_2:(a,b)\mapsto(-b,6+a).
\]
The union of these loci has finite image in the completed orbifold base and
is absent from a punctured cusp neighborhood.  Consequently this particular
elliptic/Landau sector does not produce the cusp soft modes.  This statement
does not classify all elliptic quotients, all N\'eron--Severi classes, or all
smooth integral magnetic fluxes.  See
\cref{cusp-spec-first-coordinate-descent,%
cusp-spec-elliptic-descent-finiteness,%
cusp-spec-elliptic-descent-scope}.

\item \Verified\ The fibrewise small eigenvalues are assembled into one
specified gapless operator.  On
\[
 \mathscr H_{\mathrm{aux}}
 =L^2\bigl([T_0,\infty),dT;
 L^2_0(\mathbb R^4/\mathbb Z^4,dy)\bigr)
 \cong\bigoplus_{k\in\mathbb Z^4\setminus\{0\}}
 L^2([T_0,\infty),dT),
\]
put
\[
 \lambda_k(T)=4\pi^2\lvert P_T^{-\mathsf T}k\rvert^2,\qquad
 (\mathscr L_{\mathrm{aux}}u)_k(T)=\lambda_k(T)u_k(T),
\]
\[
 D(\mathscr L_{\mathrm{aux}})=\left\{u:
 \sum_{k\ne0}\int_{T_0}^{\infty}
 \lambda_k(T)^2\lvert u_k(T)\rvert^2\,dT<\infty\right\}.
\]
This multiplication operator is nonnegative and self-adjoint,
\(\ker\mathscr L_{\mathrm{aux}}=0\), and the escaping-support Weyl
sequence in the \(k=(1,0,0,0)\) channel proves
\[
 0\in\sigma_{\mathrm{ess}}(\mathscr L_{\mathrm{aux}}).
\]
This is a spectral-gap theorem for the displayed auxiliary direct integral,
not for the complex threefold without a metric and not for a quantum gauge
Hamiltonian; see \cref{cusp-interpretation-direct-integral-gapless}.

\item \Verified\ The cover index and the cubical mesh are independent.
For integers \(M_{\mathrm{cov}}\ge1\) and
\(N_{\mathrm{mesh}}\ge2\), put
\[
 S_{M_{\mathrm{cov}}}=\operatorname{diag}(1,1,M_{\mathrm{cov}},
 M_{\mathrm{cov}}),\qquad
 P_{T,M_{\mathrm{cov}}}=P_TS_{M_{\mathrm{cov}}},
\]
\[
 \widetilde F_{T,M_{\mathrm{cov}}}
 =\mathbb R^4/P_{T,M_{\mathrm{cov}}}\mathbb Z^4
 \longrightarrow F_T,\qquad
 \operatorname{Deck}\cong(\mathbb Z/M_{\mathrm{cov}}\mathbb Z)^2.
\]
This cover has degree \(M_{\mathrm{cov}}^2\), Gram matrix and covolume
\[
 G_{T,M_{\mathrm{cov}}}
 =P_{T,M_{\mathrm{cov}}}^{\mathsf T}P_{T,M_{\mathrm{cov}}},
 \qquad
 \mathcal D_{T,M_{\mathrm{cov}}}=M_{\mathrm{cov}}^2D_T.
\]
For the cubical complex write, once and for all,
\[
 V_{N_{\mathrm{mesh}}}
   =(\mathbb Z/N_{\mathrm{mesh}}\mathbb Z)^4,\qquad
 E_{N_{\mathrm{mesh}}}
   =V_{N_{\mathrm{mesh}}}\times\{1,2,3,4\}.
\]
The diffeomorphism
\[
 \Phi_{T,M_{\mathrm{cov}}}:\mathbb R^4/\mathbb Z^4
 \xrightarrow{\ \sim\ }\widetilde F_{T,M_{\mathrm{cov}}},
 \qquad[y]\longmapsto[P_{T,M_{\mathrm{cov}}}y]
\]
transports the whole standard \(N_{\mathrm{mesh}}\)-cubical complex; its
vertex map is precisely
\[
 (\mathbb Z/N_{\mathrm{mesh}}\mathbb Z)^4
 \xrightarrow{\ \sim\ }
 \widetilde F_{T,M_{\mathrm{cov}}}[N_{\mathrm{mesh}}],
 \qquad[n]\longmapsto
 \left[\frac{P_{T,M_{\mathrm{cov}}}n}{N_{\mathrm{mesh}}}\right],
\]
so the finite set is the torsion subgroup, not the entire torus.  Its
oriented edge is
\[
 \gamma_{n,j}(t)=
 \left[\frac{P_{T,M_{\mathrm{cov}}}(n+t\mathbf e_j)}
 {N_{\mathrm{mesh}}}\right],\qquad0\le t\le1.
\]
For a compact matrix Lie group \(G\), the parallel-transport convention
\[
 U'(t)=-A_{\gamma(t)}(\dot\gamma(t))U(t),\qquad U(0)=1
\]
defines the typed map
\[
 \operatorname{Hol}_{T,M_{\mathrm{cov}},N_{\mathrm{mesh}}}:
 \Omega^1(\widetilde F_{T,M_{\mathrm{cov}}};\mathfrak g)
 \longrightarrow G^{E_{N_{\mathrm{mesh}}}},
 \qquad A\longmapsto(U_A(\gamma_{n,j}))_{n,j},
\]
and for \(A^g=g^{-1}Ag+g^{-1}dg\) its exact endpoint law is
\[
 \operatorname{Hol}(A^g)_{n,j}
 =g(\gamma_{n,j}(1))^{-1}\operatorname{Hol}(A)_{n,j}
  g(\gamma_{n,j}(0)).
\]
See \cref{cusp-link-cover-metric,cusp-link-cubical-realization,%
cusp-link-holonomy-map}.

\item \Verified\ The compact-link \(\mathrm{SU}(2)\) flux uses the
inverse-transport convention and has an exact metric-weighted action.  With
\[
 H_{\mathrm{link}}=\begin{pmatrix}i&0\\0&-i\end{pmatrix},\qquad
 \lVert H_{\mathrm{link}}\rVert^2=1,\qquad
 \exp(2\pi H_{\mathrm{link}})=1,\qquad
 N=N_{\mathrm{mesh}},
\]
set
\[
\begin{aligned}
 U_1(n)&=\exp(-2\pi n_2H_{\mathrm{link}}/N^2),\\
 U_2(n)&=
 \begin{cases}
  1,&0\le n_2\le N-2,\\
  \exp(2\pi n_1H_{\mathrm{link}}/N),&n_2=N-1,
 \end{cases}\\
 U_3(n)&=U_4(n)=1.
\end{aligned}
\]
For
\[
 U_{ij}(n)=U_i(n)U_j(n+\mathbf e_i)
 U_i(n+\mathbf e_j)^{-1}U_j(n)^{-1},
\]
including all wrap-around edges,
\begin{align*}
 U_{12}(n)&=\exp(2\pi H_{\mathrm{link}}/N^2),\\
 U_{ij}(n)&=1\quad(i<j,\ (i,j)\ne(1,2)).
\end{align*}
The principal logarithm is unambiguous because
\(0<2\pi/N^2\le\pi/2\), and
\[
 \mathcal F_N^{\log}(n)
 =N^2\operatorname{Log}(U_{12}(n))\,dy^1\wedge dy^2
 =2\pi H_{\mathrm{link}}\,dy^1\wedge dy^2.
\]
With the full inverse Gram metric,
\[
 S^{\log}_{T,M_{\mathrm{cov}},N}(U)
 =\frac{\mathcal D_{T,M_{\mathrm{cov}}}}{2g_{\mathrm{YM}}^2N^4}
 \sum_{n\in V_{N_{\mathrm{mesh}}}}
 \lVert\mathcal F_N^{\log}(n)\rVert^2_{
 G_{T,M_{\mathrm{cov}}}^{-1}\otimes\mathfrak{su}(2)}
 =\frac{2\pi^2M_{\mathrm{cov}}^2}{g_{\mathrm{YM}}^2D_T}.
\]
There is no canonical scalar Wilson action for a general nonorthogonal
metric obtained by independent scalar plaquette weights.  For this
configuration the explicitly defined \((1,2)\)-sector functional has
\[
 \kappa^{12}_{T,M_{\mathrm{cov}}}
 =(G_{T,M_{\mathrm{cov}}}^{-1})_{11}
  (G_{T,M_{\mathrm{cov}}}^{-1})_{22}
 -(G_{T,M_{\mathrm{cov}}}^{-1})_{12}^2
 =D_T^{-2},
\]
\[
 S^{W,12}_{T,M_{\mathrm{cov}},N}(U)
 =\frac{\mathcal D_{T,M_{\mathrm{cov}}}
        \kappa^{12}_{T,M_{\mathrm{cov}}}}{2g_{\mathrm{YM}}^2}
  \sum_{n\in V_{N_{\mathrm{mesh}}}}
  \left(2-\operatorname{Re}\operatorname{tr}U_{12}(n)\right)
 =\frac{M_{\mathrm{cov}}^2N^4}{2g_{\mathrm{YM}}^2D_T}
  \left(2-2\cos\frac{2\pi}{N^2}\right)
\]
Its exact mesh expansion is
\[
 S^{W,12}_{T,M_{\mathrm{cov}},N}(U)
 =\frac{M_{\mathrm{cov}}^2}{g_{\mathrm{YM}}^2D_T}
 \left(2\pi^2-\frac{2\pi^4}{3N^4}
       +\frac{4\pi^6}{45N^8}+O(N^{-12})\right),
\]
so its \(N\to\infty\) limit is
\(2\pi^2M_{\mathrm{cov}}^2/(g_{\mathrm{YM}}^2D_T)\).
See \cref{cusp-link-flux-theorem}.

\item \Verified\ The finite Gaussian oscillator is defined on the real
reduced Coulomb space with the physical inverse-Gram exterior metric.  For
\[
 \vartheta^j_{T,M_{\mathrm{cov}}}
 =P_{T,M_{\mathrm{cov}}}^{-\mathsf T}\mathbf e_j,\qquad
 \mathfrak d_{N_{\mathrm{mesh}}}(k)_j
 =N_{\mathrm{mesh}}
  (e^{2\pi ik_j/N_{\mathrm{mesh}}}-1),
\]
let
\[
 C^p_{N_{\mathrm{mesh}}}(T,M_{\mathrm{cov}})
 =\operatorname{Map}\!\left(
 (\mathbb Z/N_{\mathrm{mesh}}\mathbb Z)^4,
 \bigwedge^p(\mathbb R^4)^*\otimes\mathfrak{su}(2)\right).
\]
The exterior metric is
\[
 \left\langle
 \vartheta^{i_1}\wedge\cdots\wedge\vartheta^{i_p},
 \vartheta^{j_1}\wedge\cdots\wedge\vartheta^{j_p}
 \right\rangle
 =\det\left((G_{T,M_{\mathrm{cov}}}^{-1})_{i_a j_b}\right)_{a,b=1}^p,
\]
and the cochain inner product is
\[
 \langle a,b\rangle_p
 =\frac{\mathcal D_{T,M_{\mathrm{cov}}}}{N_{\mathrm{mesh}}^4}
 \sum_{n\in V_{N_{\mathrm{mesh}}}}\langle a(n),b(n)\rangle.
\]
Writing \(\mathsf T_ja(n)=a(n+\mathbf e_j)\), put
\[
 D_j^+=N_{\mathrm{mesh}}(\mathsf T_j-1),\qquad
 D_j^-a(n)=N_{\mathrm{mesh}}(a(n)-a(n-\mathbf e_j)),
\]
and then
\[
 d_N=\sum_{j=1}^4\varepsilon(\vartheta^j)D_j^+,\qquad
 d_N^*=-\sum_{j=1}^4\iota_{(\vartheta^j)^\sharp}D_j^-,
\]
\[
 \mathscr C_{T,M_{\mathrm{cov}},N_{\mathrm{mesh}}}
 =\ker(d_N^*:C^1\to C^0)\cap
  (\ker\Delta_N^{(1)})^\perp,\qquad
 \mathcal L=g_{\mathrm{YM}}^{-2}d_N^*d_N.
\]
The positive eigenvalue on the Fourier fibre
\(k\in(\mathbb Z/N_{\mathrm{mesh}}\mathbb Z)^4\setminus\{0\}\) is
\[
 \lambda_{T,M_{\mathrm{cov}},N_{\mathrm{mesh}}}(k)
 =\frac1{g_{\mathrm{YM}}^2}
 \left\lVert P_T^{-\mathsf T}S_{M_{\mathrm{cov}}}^{-1}
 \mathfrak d_{N_{\mathrm{mesh}}}(k)\right\rVert^2,
\]
with complex transverse multiplicity \(9\), and
\[
 \lambda(\mathbf e_1)
 =\frac{4N_{\mathrm{mesh}}^2\sin^2(\pi/N_{\mathrm{mesh}})}
 {g_{\mathrm{YM}}^2}
 \frac{L_T^2+m_T^2}{D_T^2},
\]
which is independent of \(M_{\mathrm{cov}}\).  On
\(L^2(\mathscr C_{T,M_{\mathrm{cov}},N_{\mathrm{mesh}}},da)\), the
Friedrichs realization of
\[
 \widehat H^{\mathrm{Gauss}}
 =-\frac12\Delta_{\mathscr C_{T,M_{\mathrm{cov}},N_{\mathrm{mesh}}}}
  +\frac12\langle a,\mathcal La\rangle
\]
has a unique Gaussian ground state and
\[
 \operatorname{gap}(\widehat H^{\mathrm{Gauss}})
 =\sqrt{\min_{k\ne0}\lambda(k)}
 \le\frac{2N_{\mathrm{mesh}}\sin(\pi/N_{\mathrm{mesh}})}
 {g_{\mathrm{YM}}}\frac{\sqrt{L_T^2+m_T^2}}{D_T}.
\]
For the explicit sequence
\[
 T_j=j^2,\qquad M_{\mathrm{cov},j}=j,\qquad
 N_{\mathrm{mesh},j}=j^3,
\]
one has simultaneously
\[
 \operatorname{sys}(\widetilde F_{T_j,M_{\mathrm{cov},j}})\to\infty,
 \quad
 \max_r\frac{\lVert P_{T_j,M_{\mathrm{cov},j}}\mathbf e_r\rVert}
 {N_{\mathrm{mesh},j}}\to0,
 \quad
 \operatorname{gap}(\widehat H^{\mathrm{Gauss}})\to0,
 \quad
 S^{\log}_{T_j,M_{\mathrm{cov},j},N_{\mathrm{mesh},j}}(U)\to0.
\]
See \cref{cusp-link-finite-spectrum,cusp-link-gaussian-theorem}.

\item \Verified\ The componentwise exponential
\[
 \operatorname{Exp}_{N_{\mathrm{mesh}}}:C^1
 \longrightarrow\mathrm{SU}(2)^{E_{N_{\mathrm{mesh}}}},
 \qquad a\longmapsto(\exp(a_e/N_{\mathrm{mesh}}))_e
\]
is only a local diffeomorphism at zero.  Its derivative is
\[
 D\operatorname{Exp}_{N_{\mathrm{mesh}}}(0)a
 =N_{\mathrm{mesh}}^{-1}a,
\]
and, if a smooth compact-link action \(S_W\) has the identity as a critical
point, then exactly
\[
 D^2(S_W\circ\operatorname{Exp}_{N_{\mathrm{mesh}}})_0[a,b]
 =N_{\mathrm{mesh}}^{-2}D^2(S_W)_{\mathbf1}[a,b].
\]
It is not globally injective: a cochain supported on one edge with value
\(2\pi N_{\mathrm{mesh}}H_{\mathrm{link}}\) has the same image as zero.  Thus
\[
 \mathcal H^{\mathrm{Gauss}}
 =L^2(\mathscr C_{T,M_{\mathrm{cov}},N_{\mathrm{mesh}}},da)
 \quad\text{and}\quad
 \mathcal H^{\mathrm{inv}}_{N_{\mathrm{mesh}}}
 =L^2(\mathrm{SU}(2)^{E_{N_{\mathrm{mesh}}}},dU)^{
   \mathrm{SU}(2)^{V_{N_{\mathrm{mesh}}}}}
\]
are not identified by a global coordinate chart or a constructed unitary
map; moreover \(\delta_U\notin L^2\) for product Haar measure.  No nonlinear
compact-link Hamiltonian, Gauss-law spectral problem, time direction,
reflection-positive transfer matrix, renormalized-scale limit, or
Osterwalder--Schrader reconstruction is supplied by these results.  See
\cref{cusp-link-exponential-local-only,cusp-link-physical-hilbert-space,%
cusp-link-dirac-not-l2}.

\item \Verified\ The reducible direct-sum connection on
\(\mathscr E_T=\mathscr L_T\oplus\mathscr L_T^{-1}\) is the non-flat
Yang--Mills critical point which, in the lifted splitting over
\(\mathbb R^4\), has the connection form
\[
 \nabla^T=\mathrm d+2\pi u_1\,\mathrm du_2\,H_{\mathrm{mag}},\qquad
 H_{\mathrm{mag}}=\begin{pmatrix}i&0\\0&-i\end{pmatrix},
\]
\[
 F_{\nabla^T}
 =\frac{2\pi}{D_T}H_{\mathrm{mag}}\,\mathrm dy_1\wedge\mathrm dy_2,
 \qquad \mathrm d_{\nabla^T}^*F_{\nabla^T}=0.
\]
The underlying principal \(SU(2)\)-bundle is smoothly trivial; after
choosing any of the proved global \(SU(2)\)-trivializations, the displayed
lifted-splitting form is carried to its gauge transform on that trivial
bundle.  It is not asserted to remain literally
\(\mathrm d+2\pi u_1\,\mathrm du_2\,H_{\mathrm{mag}}\) in an arbitrary
global trivialization.
For
\(\operatorname{YM}_T(\nabla)
=(2g_{\mathrm{YM}}^2)^{-1}\int_{F_T}|F_\nabla|^2\,
\mathrm d\operatorname{vol}_{F_T}\), its action is
\[
 \operatorname{YM}_T(\nabla^T)
 =\frac{2\pi^2}{g_{\mathrm{YM}}^2D_T}
 =\frac{2\pi^2}{g_{\mathrm{YM}}^2T^2}
 \left(1+\frac{B-2H}{T}+O(T^{-2})\right).
\]
On the background-gauge slice its classical Jacobi operator is
\[
 (\mathcal J_Ta)_\nu
 =-\sum_{\mu=1}^4\nabla_{e_\mu}\nabla_{e_\mu}a_\nu
  -2\sum_{\mu=1}^4[F_{\nu\mu},a_\mu],
\]
and an explicit smooth real unit vector satisfies
\[
 \mathrm d_{\nabla^T}^*a=0,\qquad
 \mathcal J_Ta=-\frac{4\pi}{D_T}a,\qquad
 \operatorname{YM}_T''|_{\nabla^T}(a,a)
 =-\frac{4\pi}{g_{\mathrm{YM}}^2D_T}.
\]
Thus the displayed reducible classical critical point is a saddle for every
finite \(T\); this is not a statement about a quantum Yang--Mills
Hamiltonian.  See
\cref{cusp-magnetic-su2-triviality,cusp-magnetic-ym-critical-point,%
cusp-magnetic-negative-mode}.

\item \Verified\ The monodromy-invariant full-flux class gives an exact
split-Cartan \(SU(r)\) theorem for every nonzero
\[
 \mathbf m=(m_1,\ldots,m_r)\in\mathbb Z^r,
 \qquad \sum_i m_i=0,
 \qquad s_2(\mathbf m)=\sum_i m_i^2.
\]
On the regular native fibre define
\[
 \mathscr E_{\mathbf m,T}
 =\bigoplus_{i=1}^r\mathscr L_{\eta,T}^{\otimes m_i},
 \qquad H_{\mathbf m}=i\operatorname{diag}(m_1,\ldots,m_r).
\]
The ordered splitting canonically and parallelly trivializes the determinant,
and, retaining the two opposite orientations used by the native construction,
\[
 c_1(\mathscr E_{\mathbf m,T})=0,
 \qquad
 c_2(\mathscr E_{\mathbf m,T})
 =-\frac{s_2(\mathbf m)}2[\eta_T]^2,
\]
\[
 \left\langle c_2,[F_T]_{\mathrm{mark}}\right\rangle=-6s_2(\mathbf m),
 \qquad
 k_{\mathbf m}:=
 \left\langle c_2,[F_T]_{x\text{-}\mathrm{complex}}\right\rangle
 =6s_2(\mathbf m).
\]
Here \(s_2(\mathbf m)\) is even, so this split family realizes precisely a
subfamily with \(k_{\mathbf m}\in12\mathbb Z_{>0}\).  The diagonal connection
\[
 \nabla^{\mathbf m,T}
 =\mathrm d-2\pi(6y^1\,\mathrm dy^4+y^2\,\mathrm dy^3)H_{\mathbf m},
 \qquad F_{\mathbf m,T}=-2\pi H_{\mathbf m}\eta_T
\]
is a reducible non-flat Yang--Mills critical point and has exact action
\[
 \operatorname{YM}_T(\nabla^{\mathbf m,T})
 =\frac{2\pi^2k_{\mathbf m}}{g_{\mathrm{YM}}^2}
  +\frac{\pi^2s_2(\mathbf m)S_T^2}{g_{\mathrm{YM}}^2D_T}.
\]
For the positive root \(i<j\), put \(q_{ij}=|m_i-m_j|\) and
\[
 \rho_T^*=\frac{\max\{\mathcal B_{T,+},\mathcal B_{T,-}\}}
                  {\min\{\mathcal B_{T,+},\mathcal B_{T,-}\}},
 \qquad
 N_-(T)=\#\{n\ge0:2n+1<\rho_T^*\}.
\]
Its complete negative normalized Coulomb spectrum is
\[
 \nu_{ij,n}(T)=q_{ij}
 \bigl((2n+1)\mathcal B_{T,\min}-\mathcal B_{T,\max}\bigr),
 \qquad 0\le n<N_-(T),
\]
with real multiplicity \(12q_{ij}^2\) before summing coincident numerical
bands.  Consequently
\[
 \operatorname{ind}_{\mathbb R}\mathcal J_{\mathbf m,T}
 =12rN_-(T)s_2(\mathbf m),
\]
and on the deep cusp tail, where \(N_-(T)=3\), this is
\(36r s_2(\mathbf m)=6rk_{\mathbf m}\).  The module also computes the
centralizer harmonic kernel and every odd-frequency wall nullity; negative
eigenspaces are descent directions, whereas only the zero-mode walls are
possible Lyapunov--Schmidt bifurcation parameters.  Finally, naturality of
\(c_2\) proves that this nonzero-flux bundle cannot extend over a carrier
with \(H^4(-;\mathbb Z)=0\), in particular not over a space diffeomorphic to
\(S^6\), \(S^{12}\), or \(S^{24}\).  See
\cref{cusp-native-full-flux-suN-topology,%
cusp-native-full-flux-suN-action-theorem,%
cusp-native-full-flux-suN-Morse-theorem,%
cusp-native-full-flux-suN-nullity-theorem,%
cusp-native-full-flux-suN-stability-boundary}.

\item \Verified\ The Koide relation has an exact ambient
quadratic-form interpretation, and it is rigorously separated from the
regular canonical cusp family.  For
\[
 Q_K(r)=\frac23(r_1+r_2+r_3)^2-(r_1^2+r_2^2+r_3^2),\qquad
 \begin{pmatrix}h\\x\\y\end{pmatrix}
 =\begin{pmatrix}
 1/\sqrt3&1/\sqrt3&1/\sqrt3\\
 1/\sqrt2&0&-1/\sqrt2\\
 1/\sqrt6&-2/\sqrt6&1/\sqrt6
 \end{pmatrix}
 \begin{pmatrix}r_1\\r_2\\r_3\end{pmatrix},
\]
one has \(Q_K(r)=h^2-x^2-y^2\).  The real-linear coefficient-space
isomorphism
\[
 (h,x,y)\longmapsto(t,m,\varrho)
 =\left(h+y,\frac{x}{\sqrt6},h-y\right)
\]
satisfies, with
\[
 (t,m,\varrho)
 =(\operatorname{Im}\tau,\operatorname{Im}\mu,
   \operatorname{Im}\beta),
\]
for the canonical realification of
\(Z=\bigl(\begin{smallmatrix}6\mu&\tau\\\beta&\mu\end{smallmatrix}\bigr)\),
\[
 \det P(\tau,\mu,\beta)
 =t\varrho-6m^2
 =\det\begin{pmatrix}t&\sqrt6\,m\\\sqrt6\,m&\varrho\end{pmatrix}
 =Q_K(r).
\]
For \(r\in\mathbb R_{\ge0}^3\setminus\{0\}\), the Koide equation
\(Q_K(r)=0\) makes this symmetric matrix positive semidefinite of rank
one and makes \(P\) singular.  Its columns therefore do not form a full
lattice and do not define a regular compact complex two-torus.  In contrast,
the actual positive radial cusp has
\[
 \det P_T=-D_T<0,\qquad
 D_T=T^2+(2H-B)T+O(1),\qquad
 \det P_T\sim-T^2.
\]
Thus the result is an isomorphism of ambient coefficient spaces with their
quadratic discriminants, not a map from Koide triples to fibres of the
constructed \(S^6\) family; see
\cref{koide-s6-thm-ambient-discriminant,koide-s6-cor-ambient-only,%
koide-s6-cor-actual-cusp}.

\item \Verified\ Independently of that discriminant comparison, the positive
defect-one shell has an exact integral parametrization and a power-saving
count with its exact constant.  Put
\[
 \Delta(x,y,z)=4(xy+xz+yz)-(x^2+y^2+z^2)
\]
and let \(R(X)\) count triples
\(0<x\le y\le z\le X\) with \(\Delta=-1\) on the terminal Vieta wall
\(z\ge4(x+y)\).  The integral change of variables
\[
 n=z-4(x+y),\qquad q=y-x,\qquad a=x+y-4n
\]
is a bijection onto
\[
 a^2-3q^2-18n^2=-2,\qquad
 a>0,\ q,n\ge0,\ a+4n>q,\qquad z=4a+17n.
\]
Write
\[
 \mathscr X=\{(a,q,n)\in\mathbb A^3:
                 a^2-3q^2-18n^2=-2\}.
\]
The exact local densities are
\[
 \sigma_2=\frac32,\qquad \sigma_3=2,\qquad
 \sigma_p=\frac{1-p^{-2}}{1-(\frac3p)p^{-1}}\qquad(p\ge5),
\]
and the Poincar\'e-residue measure of the real chamber has leading factor
\[
 I_\infty=\frac1{2\sqrt3}\log\!\left(\frac3{2\sqrt2}\right).
\]
After verifying the Gorodnik--Oh--Shah sector hypotheses and Borovoi's
adelic sign coefficient \(c_{\mathscr X}=1\), one obtains unconditionally
\[
 R(X)\sim C_0X,\qquad
 C_0=\frac9{2\pi^2}\log(2+\sqrt3)
       \log\!\left(\frac3{2\sqrt2}\right).
\]
The exact ideal parametrization now proves more.  For \(n\geq1\), put

\[
\begin{aligned}
 m_n&=18n^2-2=2(3n-1)(3n+1),\\
 \mathcal P(m_n)&=\bigl\{(C,D,E)\in\mathbb Z_{>0}^2\times\mathbb Z:\\
 &\hspace{5em}C^2D=m_n,\quad 0\leq E<D,\quad
 E^2\equiv3\pmod D\bigr\}.
\end{aligned}
\]

The map

\[
 (C,D,E)\longmapsto
 C\bigl(D\mathbb Z+(E+\sqrt3)\mathbb Z\bigr)
\]

is a bijection to the ideals of
\(\mathbb Z[\sqrt3]\) of norm \(m_n\).  Each ideal has a unique seed
\(a_0+q_0\sqrt3\) with \(a_0>2q_0\geq0\), and every solution in its orbit is

\[
 a_k+q_k\sqrt3=(a_0+q_0\sqrt3)(2+\sqrt3)^k,\qquad k\geq0.
\]

The inverse is finite and integral: apply
\((a,q)\mapsto(2a-3q,2q-a)\) while \(a\leq2q\), then recover the unique HNF
triple.  Consequently every point has multiplicity one and the number of
unit orbits at level \(n\) is

\[
 \#\mathcal P(m_n)=\rho(m_n)
 =\sum_{d\mid m_n}\chi_{12}(d)
 =\rho(3n-1)\rho(3n+1).
\]
The omitted \(n=0\) face is separate and exact:
\[
 a_k+q_k\sqrt3=(1+\sqrt3)(2+\sqrt3)^k,\qquad k\geq1,
\]
subject to the unchanged strict wall.  Its contribution is the floor count
\(R_0(X)\) of \cref{koide-param-eq:R0-exact}, and
\(R_0(X)=O(\log(2+X))\).

The exact sharp-height floor formula and its Hecke-phase Fourier expansion
are written in
\cref{koide-param-thm:integral-parametrization,koide-param-thm:exact-floor}.
A direct verification of Benoist--Oh effective well-roundedness for the
literal five-wall family, followed by their
Theorem~13.2~\cite[pp.~1931--1932]{BenoistOh_2012} on each of the
finitely many integral spin-group orbits, strengthens the asymptotic to

\[
 \boxed{R(X)=C_0X+O(X^{1-\eta})}
 \qquad\text{for some fixed }\eta>0.
\]

Neither Benoist--Oh's theorem nor its proof supplies
\(\eta\geq\tfrac12\).  The primary-source comparison is exact:
Liu--Sarnak give the smooth standard-radial exponent \(39/64\) for this
congruence arithmetic lattice, while Good and Chatzakos--Petridis give the
sharp standard-radial \(2/3\) frontier after the relevant small-spectrum
terms are displayed.  Neither theorem counts the unchanged five-wall region
with height \(4a+17n\), and neither proves its pointwise square-root
discrepancy; see \cref{koide-lit-determination}.

The exact census gives
\(R(10^8)=3{,}536{,}002\), after an independent six-coordinate replay of
all \(3{,}536{,}002\) rows.  Its largest
\(\lvert R(X)-C_0X\rvert/\sqrt X\) on
\(10^4\leq X\leq10^8\) is
\(0.1714001928088145597\ldots\), at \(X=867{,}240\); this is finite evidence
only.
The stronger proposed estimate
\(R(X)-C_0X=O_\varepsilon(X^{1/2+\varepsilon})\) is neither asserted nor
used in this status item; the proved conclusion is the displayed linear
power-saving asymptotic.  See
\cref{koide-count-thm-coordinate-bijection,koide-count-thm-local-densities,%
koide-count-thm-archimedean-factor,koide-count-thm-linear-asymptotic,%
koide-count-sec-census-1e8,koide-param-thm:power-saving,%
koide-lit-determination,%
koide-count-open-square-root-discrepancy}.

\item \Verified\ The same parity-restricted shell has an exact quaternionic
realization.  In

\[
 B=(18,-6)_{\mathbb Q}
 =\mathbb Q\langle I,J\rangle/(I^2-18,J^2+6,IJ+JI),\qquad K=IJ,
\]

put

\[
 \mathcal O=\mathbb Z1\oplus
 \mathbb Z\frac{I+J}{2}\oplus
 \mathbb Z\frac{I-J}{2}\oplus
 \mathbb Z\frac K2.
\]

The complete multiplication table proves that this is an order.  Its
reduced discriminant is \(108\), so it has index \(18\) in a maximal
overorder; the integer \(18\) is not an Eichler level, since the order is
nonmaximal at the ramified primes \(2\) and \(3\).  The algebra is a division
algebra over \(\mathbb Q\), split over \(\mathbb R\), and ramified exactly at
\(\{2,3\}\).  The bijection

\[
 (a,q,n)\longmapsto
 \gamma(a,q,n)=2+\frac{qI+aJ+nK}{2}
\]

identifies \(a^2-3q^2-18n^2=-2\), with its forced parity
\(a\equiv q\pmod2\), with the reduced-trace-\(4\), reduced-norm-\(1\) locus
of \(\mathcal O\).  Every such \(\gamma\) is primitive hyperbolic, with

\[
 \mathcal O\cap\mathbb Q(\gamma)
 =\mathbb Z[\gamma-2]\cong\mathbb Z[\sqrt3],\qquad
 C_{\mathcal O^1}(\gamma)=\{\pm\gamma^k:k\in\mathbb Z\},
\]

eigenvalues \(2\pm\sqrt3\), and translation length
\(2\log(2+\sqrt3)\).  The full shell is a finite union of
\(\mathcal O^1\)-conjugacy classes.  For the displayed real splitting, the
standard radial Huber parameter is exactly

\[
 X_i(\gamma)^2=1+\frac{(7a+15\sqrt2\,n)^2}{48},
\]

whereas the Koide height is \(4a+17n\).  Thus the hyperbolic reduction is
exact, but the radial theorem does not transfer a discrepancy estimate to
the sharp Koide polyhedron; see
\cref{koide-quat-exact-quaternion-bridge-order-shell-and-hyperbolic-conjugacy-classes,%
koide-quat-exact-reformulation-of-the-koide-chamber-count,%
koide-quat-consequence-for-the-open-square-root-discrepancy}.

\item \Verified\ The fixed-prime Erd\H{o}s--Straus divisor selector is an
exact marked slice of the candidate's terminal integral presentations.  For
\(p\equiv1\pmod{12}\) prime, put
\[
 L=\frac{p-1}{4},\qquad R_j=4j+3,\qquad
 a_j=\frac{p+R_j}{4},\qquad 0\le j<L.
\]
If \(u\in\mathbb Z_{>0}\), \(u\mid a_j^2\),
\(\varepsilon\in\{0,1\}\), and
\(R_j\mid4u+p^\varepsilon\), then
\[
 T=\frac{4u+p^\varepsilon}{R_j},\qquad
 (\ell_0,\ell_1,\ell_2)
 =\left(\frac{T+1}{4},u-jT,1\right)
\]
is integral and satisfies
\[
 12\ell_0-4\ell_1-3\ell_2=p^\varepsilon,
\]
with the exact relation-matrix identity
\[
 \begin{pmatrix}1&-j\\0&1\end{pmatrix}
 \begin{pmatrix}R_j&-u\\4&-T\end{pmatrix}
 =\begin{pmatrix}3&-\ell_1\\4&1-4\ell_0\end{pmatrix}.
\]
The inverse is \(T=4\ell_0-1\), \(u=\ell_1+jT\).  With \(p,j\), the
channel \(\varepsilon\), positive orientation, and the divisor gate retained,
this is a bijection between the complete two-channel selector and the
normalized terminal slice.  It proves an exact integral-relator crosswalk,
not a geometric morphism from the reciprocal cubic to \(X\); see
\cref{sec:es-s6-crosswalk,thm:es-s6-selector-terminal-bijection}.

\item \Verified\ The modulo-$840$ refinement has
$[N:R_{\mathrm{ES}}]=72$, $\det N=1$, and root system $A_5^4D_4$.
Theorems~\ref{esn:thm:niemeier}, \ref{esn:thm:d4-f4-triality},
\ref{esn:fa:thm:full-quotient}, and~\ref{esn:fa:thm:residue-weyl}
prove the glue, marked flag-weight comparison, full Weyl quotient and
position of the commuting residue six-cycle. For a prime
$p\equiv1\pmod{24}$ and $R>0$, $R\equiv3\pmod4$, set
$a=(p+R)/4$, $g=\gcd(R,pa)$, $R_0=R/g$, $N_0=pa/g$.
The exact divisor-to-ordered-pair map is
\[
 d>0,\quad d\mid N_0^2,\quad d\equiv-N_0\pmod{R_0}
 \quad\longmapsto\quad
 (b,c)=\left(\frac{N_0+d}{R_0},
                  \frac{N_0+N_0^2/d}{R_0}\right).
\]
It is a bijection onto the positive remaining pairs solving
$4/p=1/a+1/b+1/c$; the unordered count is
$U_R(p)=(A_R(p)+\mathbf1_{R_0\mid2})/2$
(Theorem~\ref{esn:thm:shell-bijection}).
Theorem~\ref{esn:thm:sample-fourier} gives
$Q_{R,P}=6W_{R,P}-M_{R,P}u\in A(S_{840})\subset N$ for the
retained finite sample $P$ of primes reducing into
$S_{840}=\langle289\rangle$. This is not a residue-only existence test.

\item \Verified\ The two-column \(SU(3)\) Gram quotient is the
four-dimensional positive Hermitian cone:
\[
 \operatorname{Mat}_{3\times2}(\mathbb C)/SU(3)
 \cong\operatorname{Herm}_2(\mathbb C)_{\geq0}.
\]
Its one-point compactified representation sphere satisfies
\[
 S^{12}_{\rm col}/SU(3)\cong B^4,
\]
with orbit-type boundary \(S^3\).  This \(S^{12}_{\rm col}\) is the column
representation sphere, not the cusp-smash \(S^{12}\); see
\cref{thm:hrgc-gram-quotient,thm:hrgc-s12-ball-quotient}.

\item \Verified\ The three-column orbit calculation produces an exact
rank-three carrier, not a conjectural dimensional analogy.  For
\(\mathscr V_3=\operatorname{Mat}_{3\times3}(\mathbb C)\),
\[
 \mathscr V_3/U(3)\cong\operatorname{Herm}_3(\mathbb C)_{\geq0}
\]
by the complete Gram invariant \(A^\dagger A\), whereas strict \(SU(3)\)
retains the alternating invariant:
\[
 \mathscr V_3/SU(3)
 \cong\{(H,z):H\geq0,\ |z|^2=\det H\}\cong\mathbb R^{10}.
\]
Thus
\[
 S^{17}/U(3)\cong B^8,\quad
 \mathscr V_3^+/U(3)\cong B^9,\quad
 S^{17}/SU(3)\cong S^9,\quad
 \mathscr V_3^+/SU(3)\cong S^{10}.
\]
The compactified rank-loss locus is a standard topologically locally flat
\(S^8\subset S^{10}\), and its complement is
\[
 S^{10}\setminus S^8\cong\mathbb R^8\times\mathbb C^*\simeq S^1,
\]
whose positive winding generator is \(\det A/|\det A|\).  The rank-one
trace-one stratum inside the rank-loss sphere is \(\mathbb CP^2\).  The
rank-three determinant is cubic; this cone is not the quadratic Lorentz cone
arising from two columns.  See
 \cref{thm:hr3c-u3-compactification,thm:hr3c-su3-invariants,%
 thm:hr3c-su-full-cone,thm:hr3c-rank-loss-pair,thm:hr3c-two-ladders}.

\item \Verified\ The selectively corrected Jordan determinant/Euler module supplies an exact
algebraic refinement, but no new geometric globalization.  In the fixed
orientation, every Hermitian matrix has a unique form
\(S+i[a]_{\times}\), and
\[
 \det\bigl(S+i[a]_{\times}\bigr)=\det S-a^{\mathsf T}Sa.
\]
For every marked positive-definite symmetric form \(B\), the linear section
\[
 \ell_B(\lambda,a)=\lambda B+i[a]_{\times}
\]
has the exact positive inverse image
\[
 \ell_B^{-1}\!\left(\operatorname{Herm}_3(\mathbb C)_{\geq0}\right)
 =\{\lambda\geq0:a^{\mathsf T}Ba\leq\lambda^2\det B\},
\]
a future \((3+1)\)-Lorentz cone.  The explicit order-cone lift
\(\Theta:\operatorname{Herm}_2(\mathbb C)\to
\mathbb RI_3\oplus i\mathfrak{so}(3)\) satisfies
\[
 \det_3\Theta(H)=\frac{\operatorname{tr}H}{2}\det_2H
\]
and sends positive-semidefinite ranks \(2,1,0\) to \(3,2,0\).
For a rank-one projection \(P\), the matrix \(\Theta(P)\) has eigenvalues
\(0,1/2,1\) and is not idempotent, which separately proves that \(\Theta\)
is not a Jordan homomorphism.  Neither \(\Theta\) nor \(\ell_B\) is a
toric, period-lattice, or filling map.  The claimed index-three translation
compression is also rejected: the proved six-to-two operation is the
rank-lowering map \(P:\mathbb Z^3\twoheadrightarrow\mathbb Z^2\), with the
three even and three odd Kuhn labels mapping to the two target triangle
orbits.  See
\cref{thm:teb-hermitian-determinant-splitting,%
prop:teb-theta-order-lift,prop:teb-fan-morphism}.

\item \Verified\ The rank-three toroidal successor and its exact quotient to
the reconstructed cusp are now constructed.  The local analytic family
\[
 \pi_3:N_3\longrightarrow\Delta
\]
has smooth fibre \(E_\tau^3\), deepest local equation
\(t=z_0z_1z_2z_3\), cusp monodromy
\[
 M_3=\begin{pmatrix}I_3&I_3\\0&I_3\end{pmatrix},
\]
and filling homomorphism \(a_i\mapsto0,\ b_i\mapsto e_i\).  Its irreducible
central fibre has
\[
 (f_0,f_1,f_2,f_3)=(1,7,12,6),\qquad \chi(W_3)=6.
\]
Its normalization is the permutohedral threefold with orbit vector
\((1,14,36,24)\); the positive-codimension strata map with degrees
\(2,3,4\) onto the \(7,12,6\) quotient strata.  Hence a real
\(S^8\)-total-space fibration retaining \(W_3\) requires its remaining
singular fibres to contribute Euler sum \(-4\).
For a full-rank integral shear
\(\mathcal B:\bar L_3\hookrightarrow L_3\), the exact generalization is
\[
 \chi(W_{3,\mathcal B})
 =3!\,[L_3:\mathcal B\bar L_3]
 =3!\,|\det\mathcal B|.
\]

In the candidate decomposition
\(\Lambda_F=A_F\oplus B_F\), define
\[
 P=\begin{pmatrix}1&-1&0\\0&1&-1\end{pmatrix},
 \qquad
 Q=\begin{pmatrix}0&-1&1\\1&-1&0\end{pmatrix}.
\]
Then
\[
 \Phi=P\oplus Q:H_1(T^6;\mathbb Z)
 \twoheadrightarrow H_1(T^4;\mathbb Z)
\]
obeys the literal monodromy, filling, and period identities
\[
 \Phi M_3=M_F\Phi,\qquad c_F\Phi=Qc_3,\qquad
 PZ_3(t)=Z_2(t)Q.
\]
It extends over the toric strata to a holomorphic map of cusp families with
connected elliptic kernel.  Moreover an explicit \(J\in SL_6(\mathbb Z)\)
gives
\[
 JM_3J^{-1}\cong
 M_F\oplus\begin{pmatrix}1&1\\0&1\end{pmatrix},
\]
 and the candidate's free order-three and order-four peripheral actions lift
 on the added elliptic factor.

The latter successful lifts are distinct from the diagonal candidate
\[
 M_1=I_3\otimes S,\qquad M_2=I_3\otimes Q,\qquad M_0=I_3\otimes T,
 \qquad M_1M_2M_0=I.
\]
Although its finite factors have orders \(4\) and \(6\), the maps
\(I-M_j:T^6\to T^6\) are isogenies for \(j=1,2\), so every affine lift of
either finite action has a fixed point.

The six source tetrahedron orbits also map \(S_3\)-equivariantly to the six
marked diagonal \(A_2\) Weyl chambers: with
\(\nu_i=E_{ii}-I_3/3\), the \(\pi\)-th cell maps to the cone
\[
 \mathbb R_{\geq0}\nu_{\pi_1}
 +\mathbb R_{\geq0}(-\nu_{\pi_3})
 =\{\lambda_{\pi_1}\geq\lambda_{\pi_2}
                    \geq\lambda_{\pi_3}\}.
\]
This is an exact fan/Cartan correspondence, not a geometric identification
with the Gram determinant singularity.  The global period-and-overlap gluing
over both finite-order discs remains \Open.  See
\cref{hrftd:thm:toroidal-cusp,hrftd:thm:strata,%
hrftd:cor:euler-shear-determinant,thm:teb-holomorphic-cusp-bridge,%
 prop:teb-explicit-conjugacy,prop:teb-finite-monodromy-lift,%
 prop:teb-weyl-chambers}.

\item \Verified\ The \(\mathbb Z^6\) native to the higher-rung cusp is
\[
 \Lambda_{\rm tor}=H_1(T^6;\mathbb Z)
 =L_3\oplus\overline L_3,
\]
not the cyclic group \(\mathbb Z/6\) and not the separate group
\(H^4(U;\mathbb Z)\).  For \(k=(1,1,1)\), the exact sequence
\[
 0\longrightarrow \mathbb Z(k,0)\oplus\mathbb Z(0,k)
 \longrightarrow\Lambda_{\rm tor}
 \xrightarrow{\,P\oplus Q\,}A_F\oplus B_F\longrightarrow0
\]
is \(S_3\)-equivariant and compatible with monodromy.  It splits as an
integral monodromy lattice but not as a \(\mathbb Z[S_3]\)-module.  Each
quotient copy is the \(A_2\) weight lattice; its root lattice has index three.
The six standard generators form two \(S_3\)-orbits of size three, whereas
the six Kuhn/Weyl chambers form the regular \(S_3\)-orbit.  Thus
\[
 \operatorname{rk}\Lambda_{\rm tor}=6,\qquad |S_3|=3!=6,
 \qquad\chi(W_{3,\mathcal B})=3!|\det\mathcal B|
\]
are three precisely related but nonidentical statements.  The actual map to
the Gram boundary is
\(\Phi_{\rm G}\circ\jmath\circ P\circ\operatorname{pr}_{\rm ang}\): it is
an \(S_3\)-equivariant, piecewise-linear chamber map on the real
cocharacter space, not an additive map, a toric fan morphism, or a map of
torus fibres.

The regular chamber orbit itself has an exact rank-nine realization.  With
\(\mathcal M_{\mathbb Z}=\operatorname{Hom}(\overline L_3,L_3)\),
\[
 \mathcal M_{\mathbb Z}
 =\bigoplus_i\mathbb ZE_{ii}\oplus
  \bigoplus_{i\ne j}\mathbb ZE_{ij},
\]
the second summand is \(\mathbb Z[S_3]\), and
\(c_{(\pi_1\pi_2\pi_3)}\mapsto E_{\pi_1,\pi_3}\) identifies its basis
with the six chambers.  These are precisely the six mixed basis pairs with
zero alternating pairing; the three diagonal matrix units encode the
excluded pairs with pairing one.  The real matrix carrier is
\(S_3\)-equivariantly isomorphic to \(\operatorname{Herm}_3(\mathbb C)\),
but the integral symmetric/skew splitting has index \(8\), and every
off-diagonal matrix unit maps to an indefinite Hermitian matrix.  Hence this
is an exact representation and incidence theorem inside the
nine-dimensional carrier, not a positive-semidefinite or toroidal map.  See
\cref{rsl:thm:s3-sequence,rsl:thm:three-sixes,rsl:prop:gram-scope,%
rnmt:thm:regular-torsor,rnmt:prop:representation-bridge}.

\item \Verified\ The marked angular epimorphism now gives an exact positive
observation morphism from the full rank-nine carrier to Minkowski space,
while retaining and classifying every fibre direction.  Put
\[
 P=\begin{pmatrix}1&-1&0\\0&1&-1\end{pmatrix},\quad
 u=(1,1,1)^{\mathsf T},\quad C=PP^*,\quad
 R=P^*C^{-1},\quad \Pi=uu^*/3,\quad X=(R\;u).
\]
Then \(PR=I_2\), \(RP=I_3-\Pi\), and \(\det X=1\).  The maps
\[
 \mathcal O_P(H)=PHP^*,\qquad \jmath_P(G)=RGR^*
\]
satisfy
\[
 \mathcal O_P\jmath_P=\operatorname{id},\qquad
 \jmath_P\mathcal O_P(H)=(I_3-\Pi)H(I_3-\Pi),
\]
and every \(H\in\operatorname{Herm}_3(\mathbb C)\) has unique coordinates
\[
 H=X\begin{pmatrix}G&\psi\\\psi^*&\lambda\end{pmatrix}X^*,\qquad
 \det H=\lambda\det G-\psi^*\operatorname{adj}(G)\psi.
\]
The complete positive fibre over \(G\succeq0\) is
\[
 \psi\in\operatorname{im}G,\qquad
 \lambda\geq\psi^*G^+\psi,
\]
and the zero fibre is exactly \(\{\lambda uu^*: \lambda\geq0\}\).
For \(g\in SL_2(\mathbb C)\),
\[
 T_P(g)=RgP+\Pi=X(g\oplus1)X^{-1}
\]
acts by \((G,\psi,\lambda)\mapsto(gGg^*,g\psi,\lambda)\).  Thus
\[
 \operatorname{Herm}_3(\mathbb C)
 \cong \text{Lorentz four-vector}\oplus
 \text{defining complex spin representation}\oplus\text{real scalar}.
\]
The vector action has kernel \(\{\pm I_2\}\), but the full action is
faithful because \(-I_2\) sends \(\psi\) to \(-\psi\).  Moreover
\[
 e_2(\jmath_P(G))=\tfrac13\det G,\qquad
 \det(uu^*+\varepsilon\jmath_P(G))=\varepsilon^2\det G.
\]
Consequently \(\jmath_P\) extends properly to an embedding
\(B^4\hookrightarrow S^8\); the reverse observation on the rank-loss
boundary is nonproper and has no infinity-preserving extension.  In the
strict \(SU(3)\)-quotient its determinant phase obeys
\[
 |\omega|^2=\lambda\det G-
 \psi^*\operatorname{adj}(G)\psi.
\]
Finally, for \(Z\in\operatorname{Mat}_{3\times2}(\mathbb C)\) and retained
\(\alpha,c>0\), the exact colour-invariant Lorentz momentum and mass are
\[
 p_\alpha(Z)=\alpha Z^*Z,\qquad
 m_{\alpha,c}^2(Z)=\frac{\alpha^2}{c^2}
 \sum_{a<b}|Z_{a1}Z_{b2}-Z_{a2}Z_{b1}|^2.
\]
Rank two is equivalent to positive mass, rank at most one to zero mass, and
\(m_{\alpha,c}^2(rZ)=r^4m_{\alpha,c}^2(Z)\); hence this unrestricted
kinematic carrier is gapless.  Its exact presymplectic reduction is now
complete.  For \(m_0>0\),
\[
 (\mathcal J_2\times\mathcal H^+_{m_0,c})/
 ((x,p)\sim(x+sp,p))\cong_{\rm symp}T^*\mathcal H^+_{m_0,c},
\]
so the rank-two system is exactly the future-timelike spinless particle.  For
\(m_0=0\), the intrinsic punctured rank-one reduction is
\(T^*\mathcal N_+^\times\), but the composite multiplier action is irregular:
\[
 d(\det\circ p_\alpha)_Z=0\quad
 (\operatorname{rank}_{\mathbb C}Z\leq1),
\]
so \(\dot x=\lambda p\) does not imply \(\lambda=e\); four transverse real
rank-opening directions are first-order null without being colour gauge, and
at \(Z=0\) every \(x(\tau)\) is stationary.  On the hard
\((\lambda_0,\omega_0)\)-constraint with \(\omega_0\ne0\), the momentum map
has rank four at \(w\ne0\) and rank three at \(w=0\).  Thus its own action
forces \(\dot x=0\) above the sharp threshold and permits
\(\dot x\in\mathbb Rp\) only at the threshold.  Its determinant inequality
and multiplication/Fock spectrum are exact, but the latter is not the
Hamiltonian of that action and gives a constraint-selected free spectral gap,
not a dynamically generated interacting gap.  See
\cref{hlob:thm:split-observation,hlob:thm:block-coordinates,%
 hlob:thm:lorentz-action,hlob:thm:cone-section,hlob:cor:Y3-block-form,%
 hlob:thm:colour-momentum,wpr:thm:massive,wpr:thm:massless,%
 wpr:thm:zero,wpr:thm:packet-action}.

\item \Verified\ The determinant normal slice has now been promoted from an
exact Lorentzian analogy to a completely typed host for explicit field
models, while the additional choices and the spectral limitations remain
visible.  With
\[
 \mathcal N_P=\{Y\in\operatorname{Herm}_3(\mathbb C):Yu=0\},\qquad
 \jmath_P(G)=RGR^*,
\]
the map \(\jmath_P:\operatorname{Herm}_2(\mathbb C)\to\mathcal N_P\) is a
real-linear isomorphism with inverse \(Y\mapsto PYP^*\), and
\[
 \det(uu^*+\varepsilon\jmath_P(G))=\varepsilon^2\det G.
\]
The polarized determinant is exactly the Minkowski form
\(tt'-xx'-yy'-zz'\), so the affine slice carries a constant metric of
signature \((1,3)\), its positive-semidefinite part is the future cone, and
the lifted congruence action realizes
\[
 PSL_2(\mathbb C)\cong SO^+(1,3).
\]
The standard Hermitian metrics used to define the Moore--Penrose section
select
\[
 W_P=R^*R=C^{-1},\qquad
 K_P=\{a\in SL_2(\mathbb C):a^*W_Pa=W_P\}\cong SU(2),
\]
and the positive invariant form
\(E_P(G,K)=\operatorname{tr}(GW_PKW_P)\).  The real normal representation is
\(\boldsymbol1\oplus\boldsymbol3\); its centre acts trivially.  Promoting
\(K_P\) to a local gauge symmetry gives the explicit nonnegative action
\[
 \int_M\left(
 \frac{|F_{\mathcal A}|^2}{2g_{\rm YM}^2}
 +\frac{Z_G}{2}|D_{\mathcal A}G|_{E_P}^2
 +\frac{m_G^2}{2}|G|_{E_P}^2
 +\frac{\lambda_4}{4}(\det G)^2\right)d\operatorname{vol}_g,
\]
and the lift \(\widehat{\mathcal A}=R\mathcal A P\) satisfies
\(D_{\widehat{\mathcal A}}\jmath_P(G)=
\jmath_P(D_{\mathcal A}G)\).  This gauging is an additional physical
construction; the normal field cannot screen a fundamental centre charge
without the centre-nontrivial retained spinor \(\psi\) or another added
matter field.

The candidate-specific separation is now exact rather than verbal.  The
retained inverse-dual cusp matrix obeys
\[
 G_{\eta,T}=\begin{pmatrix}2L_T&2m_T\\2m_T&-q_T/3\end{pmatrix},
 \qquad \det G_{\eta,T}=-\frac23D_T<0,
\]
whereas \(p_\alpha(Z)=\alpha Z^*Z\) ranges over the entire
positive-semidefinite cone.  Positive rescaling and complex congruence cannot
change this determinant sign/inertia, so a further morphism is required to
place these two matrices in one physical model.  On the Lorentzian host, the
free scalar Fock Hamiltonian has exact gap \(m_{\mathrm{KG}}\) when the
independently supplied coefficient is positive and is gapless when it is
zero.  For the massless compact gauge Gaussian, the lowest nonzero transverse
frequency is \(2\pi/L\), which closes as \(L\to\infty\).  Entangled vacua and
sectors occur in both explicitly gapped and explicitly gapless Hamiltonians.
Hence none of the determinant identity, the Gram carrier, or entanglement
alone proves a nonlinear four-dimensional mass gap.  The full-domain
comparison theorem identifies a sufficient exact bridge: if
\(\|Ru\|\geq c_R\|u\|\) and
\(\mathfrak b[Ru]\leq C_R(\mathfrak a_{\rm phys}[u]-
E_{\rm vac}\|u\|^2)\) throughout the physical excitation form domain and
\(\mathfrak b\geq\gamma\), then
\[
 \operatorname{gap}(A_{\rm phys})\geq\frac{\gamma c_R^2}{C_R}.
\]
An exact counterexample proves that an intertwiner defined on only one gapped
subspace is insufficient.  See
\cref{hnsp-normal-identification,hnsp-lorentz-theorem,hnsp-spin-theorem,%
hnsp-intrinsic-compact-representation,hnsp-intrinsic-gauge-scalar-theorem,%
hnsp-centre-screening-corollary,hnsp-candidate-spacelike,%
hnsp-free-field-gap,hnsp-gauge-gaussian-gap,hnsp-coercive-transfer}.

\item \Refuted\ A literal Fabel--Alp\"oge-style complex-fourfold
construction on a homotopy \(S^8\) cannot exist.  If
\(E=(TS^8,J)\) were complex of rank four, then
\[
 c_1(E)=c_2(E)=c_3(E)=0,\qquad c_4(E)=2u,
\]
and therefore \(p_2(TS^8)=2c_4(E)=4u\).  The eight-dimensional signature
formula, with \(p_1=0\) and signature zero, forces \(p_2=0\), a
contradiction.  A nonsingular bundle
\(T^6\to S^8\to S^2\) is also impossible: its homotopy sequence would
require a surjection \(\mathbb Z\to\mathbb Z^6\), and its Euler
characteristic would give \(2=0\).  A merely real singular fibration is not
refuted, but it must satisfy
\(\sum_i\chi(F_i)=2\) and the exact van Kampen, Wang, and integral
Mayer--Vietoris maps proved in the workbench.  For regular-piece monodromy
\(\rho\), the corrected homological Serre edge sequence is
\[
 0\longrightarrow(\textstyle\bigwedge^q\Lambda)_\rho
 \longrightarrow H_q(E;\mathbb Z)
 \longrightarrow
 H_1(B^\circ;\textstyle\bigwedge^{q-1}\Lambda)
 \longrightarrow0,
\]
whereas dually \(H^q(E;\mathbb Z)\) surjects onto the invariant group
\((\bigwedge^q\Lambda^\vee)^\rho\).  Thus every filling map must respect
the coinvariant subobject and the base-one-cycle quotient; the reversed
edge maps cannot be used in the Mayer--Vietoris test.  Every effective torus
action on \(S^8\) has rank at most four, and the sharp coordinate model is
\[
 S^8/T^4\cong\operatorname{Susp}(\Delta^3)\cong B^4.
\]
It is a real two-pole singular compactification with generic \(T^4\) orbit,
not a three-special-fibre \(T^6\)-fibration over \(S^2\).  See
\cref{hrs8:subsec:no-almost-complex,hrs8:eq:euler-fibres,%
 hrs8:eq:van-kampen-fillings,hrs8:eq:mv-tests,%
 hrs8:eq:serre-regular-part,hrs8:eq:serre-cohomology-invariants,%
 hrs8:eq:torus-rank,hrs8:eq:t4-quotient}.

\item \Verified\ The original reconstructed \(S^6\) carries a separate,
exact six-coordinate bundle problem on its regular locus.  With
\[
 D=W\sqcup S_1\sqcup S_2,
 \qquad U=X\setminus D,
\]
Alexander duality and the computed homology of the three special fibres give
\[
 H^4(U;\mathbb Z)
 \cong H_1(W;\mathbb Z)\oplus H_1(S_1;\mathbb Z)
       \oplus H_1(S_2;\mathbb Z)
 \cong\mathbb Z^2\oplus\mathbb Z^2\oplus\mathbb Z^2.
\]
Because \(U\) has CW dimension at most five, principal \(SU(3)\)-bundles
on \(U\) are classified exactly by their complete
\(c_2\in H^4(U;\mathbb Z)\).  The connecting map
\[
 \delta:H^4(U;\mathbb Z)\xrightarrow{\sim}H^5(X,U;\mathbb Z)
\]
shows that such a bundle extends across all three deleted fibres if and only
if \(c_2=0\).  This is the free lattice \(\mathbb Z^6\), not the cyclic
group \(\mathbb Z/6\mathbb Z\).

\item \Verified\ The primitive fibre-degree map, its saturated kernel, the
finite normal-holonomy quotients, and the common-centre matrix maps now have a
complete theorem--proof module in
\cref{zdc:sec:bridge,zdc:thm:fibre-degree,zdc:prop:finite-reductions,%
zdc:thm:tensor-centres,zdc:thm:common-centre}.  The result is exact but stops
before any gauge-bundle, connection, confinement, or spectral assertion.

In the marked basis
\[
 (a_0,b_0,a_1,b_1,a_2,b_2)
 =(\bar\gamma,\bar u,g_1,h_1,g_2,h_2),
\]
solving the transpose of the verified one-shot collar matrix gives
\[
 \operatorname{res}_F\!\left(\sum_{i=0}^2(n_i a_i+r_i b_i)\right)
 =(12n_0+4n_1-3n_2)\omega_F,
 \qquad [\deg_F]=(12,0,4,0,-3,0),
\]
and
\[
 \ker\deg_F
 =\mathbb Z\langle
 b_0,b_1,b_2,a_0-3a_1,3a_1+4a_2\rangle.
\]
The map is primitive, so restriction to the oriented fibre is surjective.
Define
\[
 q_{\rm col}:L\longrightarrow H_1(X^\circ;\mathbb Z),\qquad
 a_0\mapsto c,\quad a_1\mapsto\bar g_1,\quad a_2\mapsto\bar g_2,\quad
 b_0,b_1,b_2\mapsto0.
\]
If \(\Phi_1:H_1(X^\circ;\mathbb Z)\to\mathbb Z\) has row
\((12,4\ell_1,3\ell_2)\) in the basis \((c,\bar g_1,\bar g_2)\), then the
factorization through the actual collar quotient is
\[
 \deg_F=\Phi_1\circ q_{\rm col},\qquad
 \Phi_1(j_*[\widetilde a_0])=p.
\]

Let \(\vartheta_1(a_1)=e^{-2\pi i/3}\) and
\(\vartheta_2(a_2)=e^{-2\pi i/4}\) be the normal-line holonomy characters,
trivial on the other marked generators.  The candidate twists select
\[
 \chi_{12}\!\left(\sum_{i=0}^2(n_i a_i+r_i b_i)\right)
 =\vartheta_1(n_1a_1+r_1b_1)
  \vartheta_2(n_2a_2+r_2b_2)^{-1}
 =e^{-2\pi i(4n_1-3n_2)/12}.
\]
For \(\rho_m(k)=e^{-2\pi ik/m}\), the exact typed identities are
\[
 \boxed{\chi_{12}=\rho_{12}\circ\deg_F,\qquad
        \chi_6:=\chi_{12}^2=\rho_6\circ\deg_F.}
\]
Consequently
\[
 H^4(U;\mathbb Z)/\ker\chi_{12}\cong\mathbb Z/12\mathbb Z,
 \qquad
 H^4(U;\mathbb Z)/\ker\chi_6\cong\mathbb Z/6\mathbb Z,
\]
whereas ordinary coefficient reduction gives
\[
 H^4(U;\mathbb Z/6)\cong(\mathbb Z/6)^6.
\]
Thus this six-coordinate coefficient group is not the single distinguished
cyclic quotient selected by \(\chi_6\).  A principal \(SU(3)\)-bundle
\(P\to U\) satisfies
\[
 \chi_6(c_2(P))
 =\exp\!\left(-\frac{2\pi i}{6}
 \langle c_2(P)|_F,[F]\rangle\right).
\]
The order-six quotient is therefore the reduction of a geometrically
computed degree and of the combined local normal holonomy, not a torsion
subgroup inferred from the rank-six notation.

Tensor product gives an independent group-theoretic match:
\[
 \iota:SU(2)\times SU(3)\hookrightarrow SU(6),\qquad
 (A,B)\longmapsto A\otimes B,
 \qquad
 \mu_2\times\mu_3\xrightarrow{\sim}Z(SU(6))=\mu_6.
\]
Since \(U\) has homotopy dimension at most five, \(c_2\) classifies
\(SU(n)\)-bundles for \(n\ge3\).  For determinant-trivial bundles of ranks
two and three,
\[
 c_2(E_2\otimes E_3)=3c_2(E_2)+2c_2(E_3).
\]
Thus every chosen \(SU(6)\)-bundle class
\(h\in H^4(U;\mathbb Z)\) admits a tensor-subgroup reduction, with all
\(c_2\)-factorizations
\[
 (e,f)=(h+2t,-h-3t),\qquad t\in H^4(U;\mathbb Z).
\]
Neither \(h\) nor a reduction is preferred.  The separate support group is
\[
 H^5(U;\mathbb Z)
 \cong\ker\{\mathbb Z^3\xrightarrow{(1,1,1)}\mathbb Z\}
 \cong Q(A_2),
\]
for the combinatorial form making the three distinguished support basis
vectors orthonormal; this is not a cup-product or intersection form.  It is
not an intrinsic \(A_2\otimes\mathbb Z^2\) inside \(H^4(U)\), and the
rank-five kernel of \(\deg_F\) is not thereby an \(A_5\) root lattice.

The finite-group isomorphism
\(H^4(U)/\ker\chi_6\cong Z(SU(6))\) supplies neither transition functions
nor a connection.  The injective tensor map does not factor through
\((SU(2)\times SU(3))/\mu_6\), because its central \(\mu_6\) maps to
nontrivial scalar matrices.  The separate common-centre construction is
\[
\begin{aligned}
 \Gamma_6
 &=\bigl\{((\lambda\mu)^{-1},\lambda I_2,\mu I_3):
             \lambda\in\mu_2,\ \mu\in\mu_3\bigr\},\\
 \frac{U(1)\times SU(2)\times SU(3)}{\Gamma_6}
 &\lhook\joinrel\longrightarrow U(6),\\
 [z,A,B]&\longmapsto z(A\otimes B),
\end{aligned}
\]
where \(\Gamma_6\cong\mu_6\).  Its determinant is \(z^6\), and its image
meets \(SU(6)\) in exactly \(\iota(SU(2)\times SU(3))\); the image is not
generally contained in \(SU(6)\).  No present morphism identifies this degree-four quotient
with the degree-two lifting obstruction of a common-centre gauge quotient.
Global \(SU(3)\)-bundles on \(S^6\) are instead classified by
\(k\in\pi_5(SU(3))\cong\mathbb Z\), with \(c_3=\pm2ku\), and all restrict
trivially to \(U\).  See
\cref{hrs8:eq:h4-regular-locus,hrs8:eq:su3-classification,%
hrs8:eq:relative-delta,hrs8:eq:extension-criterion,%
hrs8:eq:c3-clutching}.

\item \Verified\ The exact operator-theoretic bridge required for any
physical no-gap inference is proved as a typed implication.  Let \(H\ge0\)
be self-adjoint on \(\mathcal H_{\mathrm{phys}}\), put
\[
 P_0=\mathbf1_{\{0\}}(H),\qquad
 \mathcal H_{\mathrm{exc}}=(I-P_0)\mathcal H_{\mathrm{phys}},\qquad
 \mathfrak h[\psi]=\lVert H^{1/2}\psi\rVert^2,\qquad
 \mathcal Q(H)=D(H^{1/2}).
\]
For a nonnegative quadratic form \(\mathfrak q_n\) and nonzero
\(u_n\in D(\mathfrak q_n)\subset\mathcal K_n\), and for
\(J_n:\operatorname{span}\{u_n\}\to\mathcal Q(H)\), define
\(\psi_n=(I-P_0)J_nu_n\).  If, for some \(c>0\),
\(C_n\ge0\), and \(\varepsilon_n\ge0\),
\[
 \lVert\psi_n\rVert\ge c\lVert u_n\rVert,\qquad
 \mathfrak h[\psi_n]\le C_n\mathfrak q_n[u_n]
 +\varepsilon_n\lVert u_n\rVert^2,
\]
\[
 C_n\frac{\mathfrak q_n[u_n]}{\lVert u_n\rVert^2}
 +\varepsilon_n\longrightarrow0,
\]
then the self-adjoint restriction
\[
 H_{\mathrm{exc}}
 =H\big|_{D(H)\cap\mathcal H_{\mathrm{exc}}}
\]
satisfies
\[
 \inf\sigma(H_{\mathrm{exc}})=0,\qquad
 0\in\sigma_{\mathrm{ess}}(H_{\mathrm{exc}}).
\]
The full spectral projection \(P_0\), rather than orthogonality to one
chosen vacuum, is essential unless vacuum uniqueness is independently
proved; see \cref{cusp-interpretation-typed-physical-transfer}.

\item \Open\ The hypotheses needed to apply that bridge to
four-dimensional quantum Yang--Mills theory have not been constructed here.
In particular, the workbench has no common gauge-invariant positive-metric
physical Hilbert space \(\mathcal H_{\mathrm{phys}}\), no reconstructed quantum
Hamiltonian \(H\), and no maps \(J_n\) satisfying the displayed norm and
energy estimates with control through gauge reduction, reflection-positive
reconstruction, renormalization, and the continuum and infinite-volume
 limits.  The exact radial scalar, continuum Hessian, compact-link flux,
 finite cochain, Gaussian-oscillator, magnetic, and split-Cartan
 \(SU(r)\) Jacobi calculations are
\Verified\ auxiliary geometric, finite-dimensional, or classical statements.
They neither prove nor disprove a mass gap for a
physical four-dimensional quantum Yang--Mills Hamiltonian; see
\cref{cusp-gauge-missing-physical-bridge,%
cusp-interpretation-present-boundary}.

\end{enumerate}

\paragraph{External validation note.}
A further context-isolated AI derivation independently recovered the cyclic
invariant lattices (11.4), determinant divisor (11.6), finite residue
cancellation (11.17), scalar test (11.19)--(11.20), all three absolute
cotangent direct images, and \(h^{1,2}=1\), \(h^{1,1}=2\).
Independent human specialist replication has not been recorded.  This is an
external-validation boundary, not a mathematical premise or an unresolved
step of the Hodge theorem; see \cref{record:h12-direct-images}.

\paragraph{Independent review checkpoint.}
Reviewer~B independently examined the reconstruction existing up to that
report's evidentiary cutoff and agreed that the intrinsic nearby-cycle,
Leray, and based van Kampen route does not depend on the then-open collapse
route.
Its disposition was provisionally favourable with focused verification
requests.  The later Codex proofs recorded here are separate work: they are
not retroactively attributed to Reviewer~B, and its report remains useful as
an independent comparison against the argument it actually saw.  Reviewer~B's
26 August analytic follow-up also independently proves the uniform
top-degree base-change corollary and the conductor--duality nonvanishing from
CAND-X, CAND-W, and TOP-H2 \cite{ReviewerBValidation2026}; those inputs are
not certified by that follow-up itself.

\paragraph{Status language.}
\Verified\ means proved here or reproduced by an explicit certificate.
\Conditional\ means a conclusion depending on an explicitly named
established or assumed hypothesis not proved in the present argument; it
never means doubtful or politely false.  A completely proved implication
whose antecedents are part of its exact statement is \Verified.  The current
candidate-specific ledger has no \Conditional\ row.  \Open\ means that no
complete proof or counterexample is recorded.  \Refuted\ means an exact
assertion has a complete contradiction or counterexample.  The status applies
to the immediately following assertion, not to an entire paper by contagion.
